\documentclass[11pt,a4paper]{article}

\usepackage[margin=2.2cm]{geometry}
\usepackage{fontspec}
\usepackage{xcolor}
\usepackage{colortbl}
\usepackage{hyperref}
\usepackage{booktabs}
\usepackage{longtable}
\usepackage{array}
\usepackage{fancyhdr}
\usepackage{titlesec}
\usepackage{enumitem}
\usepackage{tikz}
\usepackage{tcolorbox}
\tcbuselibrary{breakable}
\usepackage{parskip}

\definecolor{bgdark}{HTML}{0F1F1A}
\definecolor{bgcard}{HTML}{1A2F26}
\definecolor{accent}{HTML}{C9A857}
\definecolor{accent2}{HTML}{E8D59E}
\definecolor{gold}{HTML}{D4AF37}
\definecolor{textlight}{HTML}{E8E0D5}
\definecolor{textmuted}{HTML}{A0988E}
\definecolor{border}{HTML}{3A5040}
\definecolor{deepgreen}{HTML}{4A8060}
\definecolor{qbg}{HTML}{0D2A1F}
\definecolor{activitybg}{HTML}{2A2015}

\setmainfont[Path=/usr/share/fonts/truetype/noto/, Extension=.ttf,
  BoldFont=NotoSans-Bold,
  ItalicFont=NotoSans-Italic,
  BoldItalicFont=NotoSans-BoldItalic]{NotoSans-Regular}
\newfontfamily\arabicfont[Path=/usr/share/fonts/truetype/noto/, Extension=.ttf,
  Script=Arabic,
  BoldFont=NotoNaskhArabic-Bold]{NotoNaskhArabic-Regular}

\makeatletter
\renewcommand{\normalcolor}{\color{textlight}}
\makeatother

\titleformat{\section}
  {\normalfont\Large\bfseries\color{accent}}
  {\thesection}{1em}{}
  [\color{border}\titlerule]
\titleformat{\subsection}
  {\normalfont\large\bfseries\color{accent2}}
  {\thesubsection}{1em}{}
\titleformat{\subsubsection}
  {\normalfont\normalsize\bfseries\color{textlight}}
  {\thesubsubsection}{1em}{}

\newtcolorbox{quranbox}{
  breakable, colback=qbg, colframe=accent, coltext=textlight,
  coltitle=accent, fonttitle=\bfseries,
  title={\small Qur'an}, boxrule=0.8pt, arc=3pt,
  left=8pt, right=8pt, top=4pt, bottom=4pt
}
\newtcolorbox{hadithbox}{
  breakable, colback=bgcard, colframe=deepgreen, coltext=textlight,
  coltitle=deepgreen, fonttitle=\bfseries,
  title={\small Hadith}, boxrule=0.8pt, arc=3pt,
  left=8pt, right=8pt, top=4pt, bottom=4pt
}
\newtcolorbox{activitybox}{
  breakable, colback=activitybg, colframe=accent2, coltext=textlight,
  coltitle=accent2, fonttitle=\bfseries,
  title={\small Family Activity}, boxrule=0.8pt, arc=3pt,
  left=8pt, right=8pt, top=4pt, bottom=4pt
}
\newtcolorbox{discussionbox}{
  breakable, colback=bgcard, colframe=accent2, coltext=textlight,
  coltitle=accent2, fonttitle=\bfseries,
  title={\small Family Discussion}, boxrule=0.8pt, arc=3pt,
  left=8pt, right=8pt, top=4pt, bottom=4pt
}
\newtcolorbox{outcomebox}{
  breakable, colback=bgcard, colframe=deepgreen, coltext=textlight,
  coltitle=deepgreen, fonttitle=\bfseries,
  title={\small Learning Objective}, boxrule=0.8pt, arc=3pt,
  left=8pt, right=8pt, top=4pt, bottom=4pt
}
\newtcolorbox{duabox}{
  breakable, colback=qbg, colframe=gold, coltext=textlight,
  coltitle=gold, fonttitle=\bfseries,
  title={\small Closing Du'a}, boxrule=0.8pt, arc=3pt,
  left=8pt, right=8pt, top=4pt, bottom=4pt
}
\newtcolorbox{quotebox}{
  breakable, colback=bgcard, colframe=border, coltext=textlight,
  boxrule=0.5pt, arc=3pt,
  left=8pt, right=8pt, top=4pt, bottom=4pt
}

\pagestyle{fancy}
\fancyhf{}
\fancyhead[L]{\color{textmuted}\small Family Halaqah Course}
\fancyhead[R]{\color{textmuted}\small From the Works of Yusuf al-Qaradawi}
\fancyfoot[C]{\color{textmuted}\small\thepage}
\renewcommand{\headrulewidth}{0.4pt}
\renewcommand{\footrulewidth}{0pt}
\renewcommand{\headrule}{\color{border}\hrule width\headwidth height\headrulewidth}

\setlist[itemize]{leftmargin=1.5em, itemsep=2pt, topsep=2pt}
\setlist[enumerate]{leftmargin=1.5em, itemsep=2pt, topsep=2pt}


\begin{document}

\pagecolor{bgdark}
\color{textlight}
\normalcolor


\begin{titlepage}
\begin{tikzpicture}[remember picture, overlay]
  \fill[bgdark] (current page.south west) rectangle (current page.north east);
\end{tikzpicture}
\thispagestyle{empty}

\vspace*{3cm}
\begin{center}
{\fontsize{28}{34}\selectfont\bfseries\color{gold} Family Halaqah Course\par}
\vspace{12pt}
{\fontsize{14}{18}\selectfont\color{accent2}\itshape A Structured 30-Week Curriculum\par}
\vspace{6pt}
{\fontsize{14}{18}\selectfont\color{accent2}\itshape from the Works of Shaykh Yusuf al-Qaradawi\par}
\vspace{24pt}
{\color{border}\rule{8cm}{0.8pt}}
\vspace{24pt}

{\fontsize{12}{16}\selectfont\color{textlight} 5 Phases \quad\textbullet\quad 30 Modules \quad\textbullet\quad Weekly Family Sessions\par}
\vspace{6pt}
{\fontsize{11}{14}\selectfont\color{textmuted} Sources: Yusuf al-Qaradawi's Books, Qur'an, and Hadith\par}
\vspace{36pt}

{\fontsize{11}{14}\selectfont\color{textlight}
\begin{tabular}{rl}
\color{accent}\textbf{Phase 1:} & \color{textlight}Roots --- Faith as the Foundation (Weeks 1--6) \\
\color{accent}\textbf{Phase 2:} & \color{textlight}Trunk --- Character and Inner Purification (Weeks 7--12) \\
\color{accent}\textbf{Phase 3:} & \color{textlight}Branches --- The Family Unit (Weeks 13--18) \\
\color{accent}\textbf{Phase 4:} & \color{textlight}Leaves --- Living Islam in the World (Weeks 19--25) \\
\color{accent}\textbf{Phase 5:} & \color{textlight}Fruit --- Trust, Hope, and Eternal Vision (Weeks 26--30) \\
\end{tabular}\par}

\vfill
{\color{border}\rule{12cm}{0.4pt}}
\vspace{8pt}
{\fontsize{10}{12}\selectfont\color{textmuted} Qaradawi Library LLM Wiki \quad\textbullet\quad June 2026\par}
\end{center}
\end{titlepage}


\section*{Family Halaqah Course --- Structured Curriculum}

\subsection*{From the Works of Shaykh Yusuf al-Qaradawi}

\textbf{Course duration::} 30 weeks - 5 phases - 6 modules per phase
\textbf{Format::} Weekly family halaqah (45–60 minutes)
\textbf{Sources::} exclusively from Yusuf al-Qaradawi's books, Qur'an, and Hadith
\textbf{Date::} 27 June 2026

\vspace{6pt}\noindent{\color{border}\rule{\textwidth}{0.5pt}}\vspace{6pt}

\subsection*{Course Architecture}

\begin{quotebox}
Phase 1: Roots --- Faith as the Foundation          (Weeks 1–6)
Phase 2: Trunk --- Character and Inner Purification  (Weeks 7–12)
Phase 3: Branches --- The Family Unit                (Weeks 13–18)
Phase 4: Leaves --- Living Islam in the World        (Weeks 19–25)
Phase 5: Fruit --- Trust, Hope, and Eternal Vision   (Weeks 26–30)
\end{quotebox}

The metaphor is a tree: faith is the root, character is the trunk, family is the branches, daily life is the leaves, and the fruit is what we offer back to Allah.

\vspace{6pt}\noindent{\color{border}\rule{\textwidth}{0.5pt}}\vspace{6pt}

\subsection*{Pedagogical Principles (from Qaradawi's own methodology)}

These principles are drawn from \textit{Education and Economy in the Sunnah} (ch-01), where Qaradawi catalogues 19 Prophetic educational guidelines:

\begin{enumerate}
\item \textbf{Gradual progression} --- "Assuming a gradual attitude in teaching and making things easy rather than difficult"
\item \textbf{Individual differences} --- "A clever teacher gives every human being his share of knowledge which suits him most"
\item \textbf{Stories and similitude} --- "Choosing the best modes (like similitude, giving examples, telling stories)"
\item \textbf{Questions and dialogue} --- "Attracting the student's attention by asking questions and leading live conversation"
\item \textbf{Mercy and sympathy} --- "Being merciful and sympathetic towards the student"
\item \textbf{Moderation} --- "Being moderate in the amount being taught and avoiding putting the student in a state of boredom"
\item \textbf{Practical application} --- "Making good use of practical situations for the sake of education and guidance"
\end{enumerate}


\vspace{6pt}\noindent{\color{border}\rule{\textwidth}{0.5pt}}\vspace{6pt}

\subsection*{Module Template}

Every weekly session follows this structure:

\begin{longtable}{|p{2cm}|p{5cm}|p{8.5cm}|}
\hline
\textbf{\color{accent}Segment} & \textbf{\color{accent}Duration} & \textbf{\color{accent}Content} \\
\hline
1. Opening & 3 min & Du'a for gathering, state intention (niyyah) \\
\hline
2. Qur'an & 5 min & Recite the verse(s) for this module (Arabic + translation) \\
\hline
3. Hadith & 3 min & Read the prophetic tradition for this module \\
\hline
4. Teaching & 15 min & Qaradawi's treatment of the topic --- key points, accessible language \\
\hline
5. Family Discussion & 15 min & 3–5 guided questions for all ages \\
\hline
6. Family Activity & 10 min & A practical action for the week ahead \\
\hline
7. Closing & 3 min & Du'a related to the topic, close with salam \\
\hline
\end{longtable}

\vspace{6pt}\noindent{\color{border}\rule{\textwidth}{0.5pt}}\vspace{6pt}

\newpage
\section{PHASE 1: ROOTS --- FAITH AS THE FOUNDATION}

\begin{quotebox}\itshape "We have honored the sons of Adam" (Qur'an 17:70). Before we talk about family rules, we must understand WHO we are --- Allah's vicegerents, honored by faith. This phase establishes the root from which everything else grows.\end{quotebox}

\subsubsection*{Module 1.1 --- Iman and Human Dignity}
\textbf{Week 1:} 

\begin{outcomebox} Family members understand that their worth comes from Allah, not from appearance, achievement, or social status. \end{outcomebox}

\begin{quranbox} 17:70 --- "We have honored the sons of Adam; provided them with transport on land and sea; given them for sustenance things good and pure; and conferred on them special favours, above a great part of Our creation." \\ 95:4 --- "Indeed, We have created man in the best of molds." \\ 49:13 --- "Indeed, the most noble of you in the sight of Allah is the one with the most taqwa." \end{quranbox}

\begin{hadithbox} "You are the best of peoples, evolved for mankind" (Al `Imran: 110, cited in Faith and Life ch-01) \\ Divine hadith on proximity: "I am as My servant thinks I am... if he draws near to Me a hand's span, I draw near to him an arm's length" \end{hadithbox}

\paragraph*{Qaradawi's teaching (from *Faith and Life*, Chapter 1)::}  - Materialism reduces man to "a piece of flesh, blood, nerves, systems, glands, and cells" whose chemical value is a few piasters - Islam affirms: Allah created man in the best of molds, breathed His Spirit into him, commanded the angels to prostrate, made him His vicegerent on earth - The dignity of being a believer (karāmat al-mu'min) surpasses even the dignity of being human (karāmat al-insān) - Examples of Bilal ibn Rabah (a slave elevated above the masters of Quraysh) and Ribʿi ibn ʿAmir (an illiterate Bedouin confronting the Persian general)

\begin{discussionbox} 1. What makes a person valuable --- is it what they look like, what they own, or something deeper? \\ 2. How does knowing Allah honored us change the way we treat each other in this family? \\ 3. Has anyone ever made you feel small? How does this verse help? \\ 4. What does "the most noble of you is the one with the most taqwa" mean for our family? \end{discussionbox}

\begin{activitybox} Each family member writes or says one thing they appreciate about every other family member --- focusing on character, not achievement. \end{activitybox}

\begin{duabox} "\arabicfont{رَبَّنَا} \arabicfont{هَبْ} \arabicfont{لَنَا} \arabicfont{مِنْ} \arabicfont{أَزْوَاجِنَا} \arabicfont{وَذُرِّيَّاتِنَا} \arabicfont{قُرَّةَ} \arabicfont{أَعْيُنٍ}" --- "Our Lord, grant us from among our spouses and offspring comfort to our eyes." (25:74) \end{duabox}

\vspace{6pt}\noindent{\color{border}\rule{\textwidth}{0.5pt}}\vspace{6pt}

\subsubsection*{Module 1.2 --- Iman and Happiness}
\textbf{Week 2:} 

\begin{outcomebox} Family members understand that true happiness comes from faith, not from material possessions or entertainment. \end{outcomebox}

\begin{quranbox} 13:28 --- "Those who believe and whose hearts find rest in the remembrance of Allah. Indeed, in the remembrance of Allah do hearts find rest." \\ 12:87 --- "Do not despair of relief from Allah. Indeed, no one despairs of relief from Allah except the disbelieving people." \end{quranbox}

\begin{hadithbox} "Whoever's intention is directed to the Hereafter, Allah makes his richness in his heart and brings his self together, and the good things of worldly life come to him unwillingly." (Tirmidhi, via Anas ibn Malik, cited in Faith and Life ch-02) \end{hadithbox}

\paragraph*{Qaradawi's teaching (from *Faith and Life*, Chapter 2)::}  - Systematic evaluation: material luxury, children, science, and psychology cannot produce true happiness without faith - The Sweden case study: comprehensive social welfare (healthcare, education, generous subsidies) but widespread worry, solitude, and soaring suicide rates - American material prosperity failed to deliver contentment - True happiness flows from contentment that faith provides --- the heart oriented toward Allah receives inner richness, tranquility, and worldly provision flows naturally - Ibn al-Qayyim: spiritual suffering is the disconnection from Allah; wealth and children may themselves become divine punishment for the hypocrites

\begin{discussionbox} 1. What makes you truly happy? Is it the same as what makes you excited? \\ 2. Have you ever gotten something you really wanted and then felt empty afterward? \\ 3. What does "richness of the heart" mean? \\ 4. How can we as a family find happiness in remembrance of Allah? \end{discussionbox}

\begin{activitybox} Each family member shares one moment this week when they felt genuinely happy. Discuss: was it connected to something material, or to faith, family, or gratitude? \end{activitybox}

\begin{duabox} "\arabicfont{اللَّهُمَّ} \arabicfont{أَعِنِّي} \arabicfont{عَلَى} \arabicfont{ذِكْرِكَ} \arabicfont{وَشُكْرِكَ} \arabicfont{وَحُسْنِ} \arabicfont{عِبَادَتِكَ}" --- "O Allah, help me to remember You, thank You, and worship You in the best way." (Abu Dawud) \end{duabox}

\vspace{6pt}\noindent{\color{border}\rule{\textwidth}{0.5pt}}\vspace{6pt}

\subsubsection*{Module 1.3 --- Iman and Love}
\textbf{Week 3:} 

\begin{outcomebox} Family members understand that love --- for Allah, for family, for people --- is a fruit of faith, not separate from it. \end{outcomebox}

\begin{quranbox} 30:21 --- "And among His signs is that He created for you mates from among yourselves, that you may dwell in tranquility with them, and He has placed between you affection and mercy." \\ 67:3 --- "No want of proportion will you see in the Creation of Allah Most Gracious." \end{quranbox}

\begin{hadithbox} "By Him in Whose Hand is my soul, you will not enter Paradise unless you truly believe, and you will not truly believe unless you love one another." (Muslim) \\ "The believer does not envy or harbor hatred against anyone." (cited in Faith and Life ch-03) \end{hadithbox}

\paragraph*{Qaradawi's teaching (from *Faith and Life*, Chapter 3)::}  - Love is the gravitational force that preserves human relations from friction and destruction --- "as the law of gravity grasps the earth and prevents planets from colliding, the law of love protects human relations from clashing" - The believer loves Allah as the source of creation, beauty, and beneficence - Love for nature: appreciating Allah's creation as an act of worship - Love for people: the believer does not envy, practices altruism (ithar), upholds tolerance as a creedal obligation - Love for calamities: the believer loves even trials because they awaken the soul's capacity for resistance

\begin{discussionbox} 1. Who do you love most? How does your faith connect to that love? \\ 2. The hadith says we cannot truly believe unless we love one another. What does that mean for our family? \\ 3. Can you love someone you envy? Why not? \\ 4. How can we show love to each other this week in a way that pleases Allah? \end{discussionbox}

\begin{activitybox} Write a family "love charter" --- 3 commitments about how family members will treat each other this week. Post it on the fridge. \end{activitybox}

\begin{duabox} "\arabicfont{اللَّهُمَّ} \arabicfont{اجْعَلْ} \arabicfont{فِي} \arabicfont{قَلْبِي} \arabicfont{نُورًا}" --- "O Allah, place light in my heart." (part of the comprehensive du'a from Bukhari) \end{duabox}

\vspace{6pt}\noindent{\color{border}\rule{\textwidth}{0.5pt}}\vspace{6pt}

\subsubsection*{Module 1.4 --- Iman and Hope}
\textbf{Week 4:} 

\begin{outcomebox} Family members understand that hope is a source of strength, and that despair is contrary to faith. \end{outcomebox}

\begin{quranbox} 12:87 --- "Do not despair of relief from Allah. Indeed, no one despairs of relief from Allah except the disbelieving people." \\ 15:56 --- "Who despairs of the mercy of his Lord except those astray?" \\ 65:2-3 --- "Whoever fears Allah --- He will make for him a way out. And will provide for him from where he does not expect." \end{quranbox}

\begin{hadithbox} "Oh Allah, there is no omen but Your omen, and no good but Your good. There is no god but You." (Ibn Wahb, cited in Ethics in Islam section-04-02) \\ "Whoever changes his plans because of an omen, he has associated with Allah." (Ahmed) \end{hadithbox}

\paragraph*{Qaradawi's teaching (from *Faith and Life*, Chapter 4)::}  - Hope is "the elixir of life, motivator of its energy, alleviator of its calamities, and an incentive of delight and joy therein" - Hope motivates the farmer to work hard, the merchant to travel, the student to study, the soldier to show courage, the patient to take medicine - Despair is "a poison that has a slow effect and a volcano that destroys man's activity" - Despair and disbelief are interrelated --- each is a cause and result of the other - The story of Prophet Ayyub (Job): lost friends, children, health, wealth; wife sold her hair for sustenance; never lost faith; Allah restored everything - Ibn Mas'ud: "Loss is involved in two things: despair and self-satisfaction" - Al-Ghazali: "Happiness is only obtained through persistent pursuit and continuous perseverance"

\begin{discussionbox} 1. What is something you hoped for and got? What is something you are hoping for now? \\ 2. How does the story of Prophet Ayyub inspire our family? \\ 3. What is the difference between hope and wishful thinking? \\ 4. When something bad happens, how can our faith help us stay hopeful? \end{discussionbox}

\begin{activitybox} Create a "family hope jar" --- each member writes one hope or goal for the next month. Read them together at the end of the month and discuss how Allah answered them. \end{activitybox}

\begin{duabox} "\arabicfont{اللَّهُمَّ} \arabicfont{إِنِّي} \arabicfont{أَعُوذُ} \arabicfont{بِكَ} \arabicfont{مِنَ} \arabicfont{الْهَمِّ} \arabicfont{وَالْحَزَنِ}" --- "O Allah, I seek refuge in You from worry and sorrow." (Bukhari, Muslim) \end{duabox}

\vspace{6pt}\noindent{\color{border}\rule{\textwidth}{0.5pt}}\vspace{6pt}

\subsubsection*{Module 1.5 --- Taqwa: God-Consciousness in Daily Life}
\textbf{Week 5:} 

\begin{outcomebox} Family members understand that taqwa is being aware of Allah in every moment, not just during worship. \end{outcomebox}

\begin{quranbox} 49:13 --- "Indeed, the most noble of you in the sight of Allah is the one with the most taqwa." \\ 2:183 --- "Fasting is prescribed for you... that you may attain taqwa." \\ 2:197 --- "And take provisions, but indeed, the best provision is taqwa." \\ 5:2 --- "Cooperate in righteousness and taqwa, but do not cooperate in sin and aggression." \end{quranbox}

\begin{hadithbox} "Truth leads to piety (taqwa), and piety leads to Paradise. A man persists in telling the truth until he is recorded with Allah as truthful." (Bukhari, Muslim, via Ibn Mas'ud) \\ "Allah does not look at your appearances and your wealth, but He looks at your hearts and your deeds." (Muslim, cited in Diversion and Arts ch-04) \end{hadithbox}

\paragraph*{Qaradawi's teaching (from *Ethics in Islam* ch-01, ch-03, ch-04; *Diversion and Arts* ch-01, ch-02)::}  - Taqwa = inner awareness of Allah's presence + outer discipline of actions - The muttaqi is one who shields himself from Allah's punishment by adhering to His commands - Taqwa balances fear (khawf) and hope (raja') --- neither despair nor complacency - Taqwa is the stated PURPOSE of fasting (2:183), pilgrimage (2:197), and sacrifice (22:37) - Taqwa in entertainment: moderate diversion is lawful, but the person of taqwa exercises self-regulation - Taqwa as sole criterion of nobility --- race, lineage, wealth count for nothing before Allah - "Good intentions, piety, or sincerity do not excuse transgression on Allah's Sovereignty" --- taqwa means submitting to Allah's boundaries even when intention seems good

\begin{discussionbox} 1. What does it mean to "be aware of Allah" when no one is watching? \\ 2. How does fasting help us build taqwa? \\ 3. If taqwa is the only thing that makes someone noble, how should we treat people who seem "less important"? \\ 4. What is one area where we as a family need more taqwa? \end{discussionbox}

\begin{activitybox} Pick one daily activity (eating, getting dressed, leaving the house) and add a conscious "Allah is watching me" awareness practice for the whole week. Share experiences next halaqah. \end{activitybox}

\begin{duabox} "\arabicfont{اللَّهُمَّ} \arabicfont{إِنِّي} \arabicfont{أَسْأَلُكَ} \arabicfont{الْهُدَى} \arabicfont{وَالتُّقَى} \arabicfont{وَالْعَفَافَ} \arabicfont{وَالْغِنَى}" --- "O Allah, I ask You for guidance, piety, chastity, and self-sufficiency." (Muslim, via Tirmidhi) \end{duabox}

\vspace{6pt}\noindent{\color{border}\rule{\textwidth}{0.5pt}}\vspace{6pt}

\subsubsection*{Module 1.6 --- Tawakkul: Trust in Allah's Plan}
\textbf{Week 6:} 

\begin{outcomebox} Family members understand that tawakkul is active trust --- tying the camel AND relying on Allah. \end{outcomebox}

\begin{quranbox} 65:3 --- "Whoever relies upon Allah --- then He is sufficient for him." \\ 3:160 --- "If Allah should aid you, no one can overcome you; but if He should forsake you, who is there that can aid you after Him? And upon Allah let the believers rely." \\ 4:81 --- "Allah is sufficient as Trustee." \end{quranbox}

\begin{hadithbox} "Tie it and rely on Allah." (Tirmidhi, via Anas ibn Malik --- the Prophet's answer to the man who asked whether to tie his camel or leave it loose) \\ The Prophet's Hijra: prepared camels, company, guide, uncommon route, cave --- then "Have no fear, for Allah is with us." (Bukhari, via Abu Bakr) \end{hadithbox}

\paragraph*{Qaradawi's teaching (from *Ethics in Islam*, section-04-01)::}  - Tawakkul is NOT passive fatalism --- it is full trust in Allah combined with taking every reasonable action - "Tie your camel and rely on Allah" --- the Prophetic model combines preparation with trust - The Prophet's Hijra demonstrates this: he prepared everything strategically, then trusted Allah completely - Those who seek refuge in Allah are never defeated; those who submit to His plans are never misguided - Tawakkul drives effort, it does not replace it

\begin{discussionbox} 1. What does "tie your camel" mean for a student? For a parent? For a child? \\ 2. Have you ever worried about something, done your best, and then Allah solved it in a way you didn't expect? \\ 3. What is the difference between trusting Allah and being lazy? \\ 4. What is one thing our family is worrying about now? How can we "tie the camel" and then trust Allah? \end{discussionbox}

\begin{activitybox} Family tawakkul practice --- identify one family challenge (school, work, health, finances). Make a concrete action plan (tying the camel). Then make du'a together asking Allah to handle what you cannot. \end{activitybox}

\begin{duabox} "\arabicfont{حَسْبُنَا} \arabicfont{اللَّهُ} \arabicfont{وَنِعْمَ} \arabicfont{الْوَكِيلُ}" --- "Allah is sufficient for us, and He is the best disposer of affairs." (Qur'an 3:173) \end{duabox}

\vspace{6pt}\noindent{\color{border}\rule{\textwidth}{0.5pt}}\vspace{6pt}

\newpage
\section{PHASE 2: TRUNK --- CHARACTER AND INNER PURIFICATION}

\begin{quotebox}\itshape "I have been sent to perfect good morals." (Hadith). Faith produces character. This phase builds the inner qualities that make a family strong.\end{quotebox}

\subsubsection*{Module 2.1 --- Akhlaq: Good Character as the Core of Islam}
\textbf{Week 7:} 

\begin{outcomebox} Family members understand that good character is not an add-on to Islam --- it IS the core. \end{outcomebox}

\begin{quranbox} 68:4 --- "And indeed, you are of a great moral character." \\ 2:177 --- "Righteousness is not that you turn your face toward the east or the west, but righteousness is [in] one who believes in Allah, the Last Day, the angels, the Book, and the prophets and gives wealth... establishes prayer and gives Zakat; [those who] fulfill their promises... and are patient in poverty and hardship." \end{quranbox}

\begin{hadithbox} "I have been sent to perfect good morals and conduct." (Ahmed 8952, al-Bukhari, al-Hakim) \\ "A person reaches the status of those who fast and pray by good character." (cited in Ethics in Islam) \end{hadithbox}

\paragraph*{Qaradawi's teaching (from *Ethics in Islam*, ch-01, ch-03, ch-04)::}  - Akhlaq = the inner self with its stable traits, good or evil - Two types: innate (gharizi) --- endowed by Allah; acquired (muktasab) --- through practice, companionship, and self-discipline - Character is MUTABLE, not fixed: "Dogs and horses can be trained to be gentle, and people too" (al-Ghazali) - Key insight from al-Ghazali: character is what you ARE, not what you APPEAR to do - Good character outweighs ritual: a person reaches the status of those who fast and pray through good character alone - Faith, ethics, and worship form an indivisible whole (Q 2:177 links creed, ritual, and ethics into a single reality) - Companionship shapes character --- applies to friends, media, environment

\begin{discussionbox} 1. What is character? Is it what you do when people are watching, or when no one is watching? \\ 2. The Prophet said he was sent to "perfect good morals." What does that tell us about what Islam is really about? \\ 3. Can a person change their character? How? \\ 4. Who in the family has a trait the others admire? What makes it admirable? \end{discussionbox}

\begin{activitybox} "Character spotlight" --- each week, one family member is spotlighted. Others share one good character trait they see in that person and one way that trait inspires them. \end{activitybox}

\begin{duabox} "\arabicfont{اللَّهُمَّ} \arabicfont{أَدِّبْنِي} \arabicfont{بِأَحْسَنِ} \arabicfont{أَدَبِكَ}" --- "O Allah, teach me the best of manners." (variant of the Prophet's du'a) \end{duabox}

\vspace{6pt}\noindent{\color{border}\rule{\textwidth}{0.5pt}}\vspace{6pt}

\subsubsection*{Module 2.2 --- Modesty (Haya'): The Distinctive Mark of Islam}
\textbf{Week 8:} 

\begin{outcomebox} Family members understand that modesty is a branch of faith that applies to everyone, and that it begins in the heart. \end{outcomebox}

\begin{quranbox} 24:30-31 --- "Tell the believing men that they should lower their gazes and guard their private parts... And tell the believing women that they should lower their gazes and guard their private parts and not display their adornment except that which appears thereof." \end{quranbox}

\begin{hadithbox} "Faith consists of 60 (or 70) branches. The best of them is to profess that 'There is no Allah but Allah'. The lowest of them is to remove harmful things from the road. Modesty is also a branch of faith." (Bukhari, Muslim) \\ "Modesty is part of faith, and it brings nothing but good." (Bukhari, Muslim, via Ibn Amr) \\ "Every religion has its distinctive characteristic, and modesty is the distinctive characteristic of Islam." (Ibn Majah, via Anas) \\ "If you have no sense of modesty, do whatever you like." (Bukhari, via Abu Masoud --- from the earlier prophets) \end{hadithbox}

\paragraph*{Qaradawi's teaching (from *Ethics in Islam*, section-04-02, pp.333-334)::}  - Modesty = being repulsed by immoral acts, rejecting anything undesirable from a religious or cultural perspective - Ibn Miskawayh: "The first sign of sharp intellect in children is modesty. It reflects their ability to discern what is good and what is not." - The Prophet was "more modest than a virgin behind her veil" (Abu Sa'id al-Khudri) - Modesty ≠ weakness, cowardice, or timidity --- "a faithful person has the courage of a lion when it comes to defending what is true and right" - Modesty should not impede learning, pursuing truth, or standing up for justice - Abu Tammam's poem: "If you do not fear the consequences of what you do at night, and if you have no sense of modesty, then do as you might. Nothing good is left in this life, if the beacon of modesty goes dark."

\begin{discussionbox} 1. What does modesty look like for a boy/man? For a girl/woman? \\ 2. How is modesty different from shyness or weakness? \\ 3. The Prophet was "more modest than a virgin" but also brave. How can someone be both? \\ 4. How does modesty apply to how we use phones, social media, and the internet? \end{discussionbox}

\begin{activitybox} Family modesty audit --- review one day's media consumption together. Ask: "Was this modest?" Discuss what changes to make. \end{activitybox}

\begin{duabox} "\arabicfont{اللَّهُمَّ} \arabicfont{اسْتُرْ} \arabicfont{عَوْرَاتِنَا} \arabicfont{وَآمِنْ} \arabicfont{رَوْعَاتِنَا}" --- "O Allah, conceal our faults and calm our fears." (variant Prophetic du'a) \end{duabox}

\vspace{6pt}\noindent{\color{border}\rule{\textwidth}{0.5pt}}\vspace{6pt}

\subsubsection*{Module 2.3 --- Humility and Honor}
\textbf{Week 9:} 

\begin{outcomebox} Family members understand that humility is strength, not weakness, and that it coexists with honor. \end{outcomebox}

\begin{quranbox} 17:24 --- "And lower to them the wing of humility out of mercy." (to parents) \\ 17:37 --- "Do not walk upon the earth exultantly. Indeed, you will never tear the earth apart, and you will never reach the mountains in height." \\ 31:18 --- "Do not turn your cheek [in contempt] toward people and do not walk through the earth exultantly. Indeed, Allah does not like the self-deluded and the boastful." \\ 63:8 --- "To Allah belongs [all] honor, and to His Messenger, and to the believers." \\ 5:54 --- "humble toward the believers, powerful against the disbelievers" \end{quranbox}

\begin{hadithbox} "He who has in his heart the weight of a mustard seed of pride shall not enter Paradise." (Muslim, via Ibn Mas'ud) \\ "Allah has revealed to me that you should humble yourselves to one another. One should not hold himself above another." (Muslim, via Iyad ibn Hammam) \\ "I am not a king, but a man whose mother ate dried meat." (Ibn Majah, via Abu Masoud) \end{hadithbox}

\paragraph*{Qaradawi's teaching (from *Ethics in Islam*, section-04-02, pp.335-336)::}  - Humility = doing away with pride and arrogance --- the traits that lead to Hell - The Prophet cleaned his own shoes, mended his clothes, milked his sheep, helped family and servants - When meeting companions, "nobody could tell which one the prophet was" - Humility and honor are NOT opposites --- the believer is "humble toward believers, powerful against disbelievers" (5:54) - Honor ≠ pride. "Whoever desires honor through power --- then to Allah belongs all honor" (35:10) - Humility is mentioned in the Qur'an specifically in the context of parents: "lower to them the wing of humility out of mercy" (17:24)

\begin{discussionbox} 1. What is the difference between being humble and having low self-esteem? \\ 2. The Prophet cleaned his own shoes and helped at home. What does that tell us about leadership in a family? \\ 3. Can you be humble and still stand up for yourself? How? \\ 4. How can we show humility to our parents and grandparents? \end{discussionbox}

\begin{activitybox} "Humble service week" --- each family member chooses one household task they normally don't do and does it quietly all week without asking for recognition. Reveal at next halaqah. \end{activitybox}

\begin{duabox} "\arabicfont{اللَّهُمَّ} \arabicfont{اجْعَلْنِي} \arabicfont{هَادِيًا} \arabicfont{مُهْتَدِيًا}" --- "O Allah, make me a guide and one who is guided." (Tirmidhi, via Ibn Abbas) \end{duabox}

\vspace{6pt}\noindent{\color{border}\rule{\textwidth}{0.5pt}}\vspace{6pt}

\subsubsection*{Module 2.4 --- Truthfulness and Trustworthiness}
\textbf{Week 10:} 

\begin{outcomebox} Family members understand that truthfulness is the bridge to Paradise and that trustworthiness is the mark of a believer. \end{outcomebox}

\begin{quranbox} 33:35 --- "Indeed, the Muslim men and Muslim women, the believing men and believing women... the truthful men and truthful women... Allah has prepared for them forgiveness and a great reward." \end{quranbox}

\begin{hadithbox} "Truth leads to piety (taqwa), and piety leads to Paradise. A man persists in telling the truth until he is recorded with Allah as truthful." (Bukhari, Muslim, via Ibn Mas'ud) \\ "The value of an action depends on the intentions behind it. A man will be rewarded only for what he intends." (Bukhari, Muslim, via Umar --- the foundational hadith on intention) \end{hadithbox}

\paragraph*{Qaradawi's teaching (from *Ethics in Islam*, ch-03, section-04-01)::}  - The chain: truth \ensuremath{\rightarrow} taqwa \ensuremath{\rightarrow} Paradise --- truthfulness is both the fruit of virtuous action AND the gateway to salvation - The Prophet was known as al-Amin (the Trustworthy) before prophethood --- character preceded message - Intention (niyyah) defines action: "Actions are but by intentions" - Sincerity (ikhlas) = purifying intention for Allah alone, free from showing off (riya') - "Good intention transforms diversion into an approach to Allah. Evil intention corrupts good deeds which are done for mere pretence rather than worship of Allah." - BUT: "Good intentions, piety, or sincerity do not excuse transgression on Allah's Sovereignty" --- intention cannot make haram into halal

\begin{discussionbox} 1. Why was the Prophet called "al-Amin" (the Trustworthy) before he was called a prophet? \\ 2. What happens when someone tells a "small" lie? Does it get easier or harder to tell the truth? \\ 3. What does it mean to do something "for Allah's sake" vs. "to be seen"? \\ 4. Can you think of a time when telling the truth was hard but you're glad you did? \end{discussionbox}

\begin{activitybox} "Trust challenge" --- each family member commits to one act of honesty this week that is difficult (admitting a mistake, returning something borrowed, telling the truth when it costs something). Share next week. \end{activitybox}

\begin{duabox} "\arabicfont{اللَّهُمَّ} \arabicfont{أَرِنِي} \arabicfont{الْحَقَّ} \arabicfont{حَقًّا} \arabicfont{وَارْزُقْنِي} \arabicfont{اتِّبَاعَهُ}" --- "O Allah, show me the truth as truth and grant me the ability to follow it." (variant du'a) \end{duabox}

\vspace{6pt}\noindent{\color{border}\rule{\textwidth}{0.5pt}}\vspace{6pt}

\subsubsection*{Module 2.5 --- Contentment and Moderation}
\textbf{Week 11:} 

\begin{outcomebox} Family members understand that contentment is the true wealth and that moderation is the nature of Islam. \end{outcomebox}

\begin{quranbox} 2:143 --- "Thus We have made you a middle nation." \\ 25:67 --- "Who when they spend, are not extravagant nor niggardly, but find a middle way between." \\ 7:31 --- "O children of Adam, take your adornment at every masjid, and eat and drink, but be not excessive. Indeed, He likes not those who commit excess." \\ 57:23 --- "In order that you not despair over what has eluded you and not exult [in pride] over what He has given you." \end{quranbox}

\begin{hadithbox} "Your body has its rights, your wife has her rights, and your guest has his rights. Give each their rights." (Bukhari 1975, Muslim 1159, via Amr ibn al-As) \\ "Richness does not lie in abundance of worldly goods, but richness is the richness of the heart." (Bukhari, Muslim, via Abu Hurayrah) \\ "Be satisfied with what Allah has allotted for you and you shall be the richest person." (Ibn Majah, via Abu Hurayrah) \\ "Whoever begins the day feeling secure with his family and healthy with enough provisions for the day is as though he owns the whole world." (Tirmidhi, via Ubaydullah ibn Muhsin) \\ "Look to those below you, and do not look to those above you." (Bukhari, Muslim, via Abu Hurayrah) \\ "Love moderately, for you may hate those you love one day. And hate moderately, for you may love those you hate one day." (Tirmidhi, via Abu Hurayrah) \end{hadithbox}

\paragraph*{Qaradawi's teaching (from *Ethics in Islam*, section-04-02, pp.340-342)::}  - Moderation is the very NATURE of Islam --- the middle-ground community - Balance in worship, family, rest, and social duties --- "excessive consumption of the permissible things can be harmful" - The Prophet: "Eat, drink, and wear [fine] clothes, as long as that does not involve any extravagance or vanity" - Moderation even in water for ablution: "Do not waste water even if you are on the bank of a flowing river" - Moderation in emotions: "Love moderately, for you may hate those you love one day" - Contentment (qana'ah) protects against greed and jealousy - "Satisfaction is never about quantity, but it is about quality. If one has all the riches in the world, it will never be enough without contentment."

\begin{discussionbox} 1. What does "richness of the heart" mean? Can you be rich without money? \\ 2. The hadith says "whoever begins the day feeling secure with his family and healthy with enough provisions is as though he owns the whole world." Do we feel that way? \\ 3. What does it mean to love moderately? Does it mean loving less? \\ 4. What is one thing we have more than enough of? What is one thing we should be grateful for? \end{discussionbox}

\begin{activitybox} "Contentment journal" --- each family member writes three things they are grateful for each day for one week. Share highlights at next halaqah. \end{activitybox}

\begin{duabox} "\arabicfont{اللَّهُمَّ} \arabicfont{قَنِّعْنِي} \arabicfont{بِمَا} \arabicfont{رَزَقْتَنِي}" --- "O Allah, make me content with what You have provided me." (variant du'a) \end{duabox}

\vspace{6pt}\noindent{\color{border}\rule{\textwidth}{0.5pt}}\vspace{6pt}

\subsubsection*{Module 2.6 --- Tazkiyah and Ikhlas: Purifying the Soul and Intention}
\textbf{Week 12:} 

\begin{outcomebox} Family members understand that spiritual purification is an active, daily process and that sincerity is the criterion for all actions. \end{outcomebox}

\begin{quranbox} 91:9-10 --- "Successful is the one who purifies it, and failed is the one who corrupts it." \\ 98:5 --- "They were not commanded except to worship Allah, sincere to Him in religion." \\ 18:110 --- "Whoever would hope for the meeting with his Lord --- let him do righteous work and not associate in the worship of his Lord anyone." \end{quranbox}

\begin{hadithbox} "Actions are but by intentions. A man will be rewarded only for what he intends." (Bukhari, Muslim, via Umar) \\ "The righteous are those who are pious and hidden. If they are absent, they are not missed, and if they are present, they are not recognized. Their hearts are beacons of guidance." (Ibn Majah, al-Hakim) \end{hadithbox}

\paragraph*{Qaradawi's teaching (from *Ethics in Islam*, section-04-01; *concept-tazkiyah*; *concept-ikhlas*)::}  - Tazkiyah = cleansing the soul from vice and cultivating virtue --- an ACTIVE, knowledge-driven process - "Moral life is founded on knowledge and willpower" --- not passive mysticism - Dual meaning from the root z-k-w: tazkiyat al-nafs (purification of SOUL) + tazkiyat al-mal (purification of WEALTH) - The Sunnah is "the agreed-upon source for the purification of the soul" - Ikhlas = purifying intention for Allah alone, free from showing off (riya') - "The value of an action depends on the intentions behind it" - People are rewarded for sincere intentions even if they don't materialize as actions - "Good intention transforms diversion into an approach to Allah" - Three divine ethics that support tazkiyah: sincerity (ikhlas), observing Allah (muraqaba), self-accountability (muhasaba) - Self-accountability: "A wise man is one who calls himself to account and does noble deeds to benefit him after death" (Tirmidhi, via Shaddad ibn Aws) - Umar: "Hold yourselves accountable before you are held accountable"

\begin{discussionbox} 1. What does "purifying the soul" mean? Is it something you do once or every day? \\ 2. How can you tell if you are doing something for Allah or for people to see? \\ 3. What is "self-accountability"? How can we practice it as a family? \\ 4. What is one bad habit we want to purify from our family life? \end{discussionbox}

\begin{activitybox} Family self-accountability practice --- each evening before bed, each family member silently reviews: (1) What did I do well today? (2) What could I have done better? (3) What will I do differently tomorrow? Share one reflection at the next halaqah. \end{activitybox}

\begin{duabox} "\arabicfont{اللَّهُمَّ} \arabicfont{طَهِّرْ} \arabicfont{قَلْبِي} \arabicfont{مِنَ} \arabicfont{النِّفَاقِ} \arabicfont{وَعَمَلِي} \arabicfont{مِنَ} \arabicfont{الرِّيَاءِ}" --- "O Allah, purify my heart from hypocrisy and my actions from showing off." (variant du'a) \end{duabox}

\vspace{6pt}\noindent{\color{border}\rule{\textwidth}{0.5pt}}\vspace{6pt}

\newpage
\section{PHASE 3: BRANCHES --- THE FAMILY UNIT}

\begin{quotebox}\itshape "And among His signs is that He created for you mates... and placed between you affection and mercy." (Qur'an 30:21). Faith and character produce strong families. This phase addresses marriage, spousal relations, parent-child bonds, and the Prophetic method of raising children.\end{quotebox}

\subsubsection*{Module 3.1 --- Marriage: The Foundation of Family}
\textbf{Week 13:} 

\begin{outcomebox} Family members understand marriage as a solemn covenant blessed by Allah, not merely a social contract. \end{outcomebox}

\begin{quranbox} 30:21 --- "And among His signs is that He created for you mates from among yourselves, that you may dwell in tranquility with them, and He has placed between you affection and mercy." \\ 4:21 --- "How could you take it when you have gone in unto each other and they have taken from you a solemn covenant?" \\ 24:32 --- "And marry such of you as are solitary and the pious of your slaves and maid-servants. If they be poor, Allah will enrich them of His bounty." \end{quranbox}

\begin{hadithbox} "I marry women too. He who turns away from my Sunnah has no relation with me." (Bukhari 5063, Muslim 1401, via Anas) \\ "There is no monasticism in Islam." (cited in Ethics in Islam ch-04) \\ "Your body has its rights, your wife has her rights." (Bukhari 1975, Muslim 1159, via Amr ibn al-As) \\ "This world is all but temporary joy, and the best temporary joy in this world is a righteous woman." (Muslim 1467, Ahmed 6567) \\ "A woman is married for four reasons: her wealth, her social status, her beauty, and her religion. Marry the religious one; otherwise you will regret it." (Bukhari 5090, Muslim 1466) \end{hadithbox}

\paragraph*{Qaradawi's teaching (from *Ethics in Islam* ch-01, ch-04, section-04-03; *Economic Security* ch-01, ch-04)::}  - Marriage is the foundation of social structure --- a covenant blessed by God (mīthāqan ghalīẓan, Q 4:21) - Islam categorically rejects celibacy: "No monasticism in Islam" --- the Prophet corrected followers who vowed extreme asceticism - Spouse selection: religion and character over wealth, status, and beauty --- "Marry the religious one; otherwise you will regret it" - Engagement (khitbah) is recognized, but couples must not be alone before the contract (khalwah) - Marriage poverty concern: the community bears responsibility to facilitate marriage for those who cannot afford it --- zakat funds may be used - The Qur'an promises: "If they be poor, Allah will enrich them of His bounty" (24:32) - Marriage creates kinship beyond the couple: "It is He who has created from water mankind and made them related through lineage and marriage" (25:54)

\begin{discussionbox} 1. What makes a marriage Islamic vs. just cultural? \\ 2. The Prophet said "no monasticism in Islam." What does that mean for how we view family life? \\ 3. Why does Islam say to choose a spouse for religion and character, not wealth or beauty? \\ 4. How can the community help young people get married? \end{discussionbox}

\begin{activitybox} Married couples share one thing they appreciate about their marriage. Children/teens discuss what qualities they would look for in a future spouse, guided by the Prophetic criteria. \end{activitybox}

\begin{duabox} "\arabicfont{اللَّهُمَّ} \arabicfont{بَارِكْ} \arabicfont{لِي} \arabicfont{فِي} \arabicfont{أَهْلِي}" --- "O Allah, bless me in my family." (variant du'a) \end{duabox}

\vspace{6pt}\noindent{\color{border}\rule{\textwidth}{0.5pt}}\vspace{6pt}

\subsubsection*{Module 3.2 --- Affection, Mercy, and Mutual Rights}
\textbf{Week 14:} 

\begin{outcomebox} Family members understand that marriage is built on affection (mawaddah) and mercy (rahmah), and that kindness is a divine command. \end{outcomebox}

\begin{quranbox} 30:21 --- "And He has placed between you affection and mercy." \\ 4:19 --- "O you who have believed, it is not lawful for you to inherit women by compulsion. And do not make difficulties for them in order to take [back] part of what you gave them unless they commit a clear immorality. And live with them in kindness, for if you dislike them --- perhaps you dislike a thing and Allah makes therein much good." \\ 2:228 --- "And women shall have rights similar to the rights against them, according to what is equitable." \end{quranbox}

\begin{hadithbox} "A believing man should not hate a believing woman; if he dislikes one of her personal traits, he will be pleased with another trait." (Muslim 1469) \\ "Each one of you is a caretaker (ra'i) and is responsible for those under his care." (Bukhari, Muslim) \\ "Wasting the sustenance of his dependants is sufficient sin for a man." (Bukhari, Muslim) \end{hadithbox}

\paragraph*{Qaradawi's teaching (from *Ethics in Islam* ch-01; *concept-marriage*)::}  - Marriage is built on mawaddah (affection) and rahmah (mercy) --- these are divine gifts, not human inventions - The Qur'an uses the second-person PLURAL pronoun in "He placed between you affection and mercy" rather than the dual --- because "the affection and mercy move from the nucleus of the married couple to the rest of their family and across society" - Kindness to wives is commanded: "Live with them in kindness" (4:19) - If you dislike something about your spouse, "perhaps you dislike a thing and Allah makes therein much good" --- the hidden blessing principle - Spouses are caretakers (ra'i) of each other and their children --- accountability before Allah - "Wasting the sustenance of his dependants is sufficient sin for a man" --- providing for family is a religious obligation, not optional generosity

\begin{discussionbox} 1. What does "affection and mercy" look like in daily married life? \\ 2. The verse says "if you dislike them, perhaps Allah makes therein much good." What does that teach us about patience in marriage? \\ 3. "Each one of you is a caretaker." What is each family member's area of responsibility? \\ 4. How can we show more affection and mercy to each other this week? \end{discussionbox}

\begin{activitybox} "Mercy challenge" --- each family member does one act of mercy for another family member this week without being asked (help with homework, cook a meal, listen to a problem, make a bed). \end{activitybox}

\begin{duabox} "\arabicfont{رَبَّنَا} \arabicfont{هَبْ} \arabicfont{لَنَا} \arabicfont{مِنْ} \arabicfont{أَزْوَاجِنَا} \arabicfont{وَذُرِّيَّاتِنَا} \arabicfont{قُرَّةَ} \arabicfont{أَعْيُنٍ}" --- "Our Lord, grant us from among our spouses and offspring comfort to our eyes." (25:74) \end{duabox}

\vspace{6pt}\noindent{\color{border}\rule{\textwidth}{0.5pt}}\vspace{6pt}

\subsubsection*{Module 3.3 --- Simplifying Weddings and Dowries}
\textbf{Week 15:} 

\begin{outcomebox} Family members understand the Prophetic principle of simplicity in marriage celebrations and the danger of cultural extravagance. \end{outcomebox}

\begin{quranbox} 4:4 --- "And give the women [upon marriage] their [bridal] gifts graciously." \\ 24:33 --- "And let those who find not the wherewithal for marriage keep themselves chaste, until Allah gives them means out of His Grace." \\ 65:7 --- "Allah does not charge a soul except [according to] what He has given it." \end{quranbox}

\begin{hadithbox} "The best of marriages are the simplest ones." / "The best of dowries are the simplest ones." (Abu Dawud 2117, al-Hakim) \\ When Ali proposed to Fatima and was asked what he possessed, he replied that he had nothing but a shield; the Prophet said: "Then, give it to her." (Abu Dawud 2125, al-Nassa'i 3375) \\ "If a man with a good moral character and religious commitment comes to propose, then marry [your daughter] to him, for if you do not do that, there will be social strife and widespread corruption." (al-Tirmidhi 1084, Ibn Majah 1967) \end{hadithbox}

\paragraph*{Qaradawi's teaching (from *Ethics in Islam* section-04-03; *Economic Security* ch-01, ch-04)::}  - Cultural dowry inflation is strongly criticized --- particularly South Asia (bride's father pays groom, reversing the Islamic mahr) and Egypt (bride's family expected to furnish home luxuriously) - These cultural patterns "make it difficult for young people to get married" and push the average marriage age upward - Excessive weddings delay marriage \ensuremath{\rightarrow} risk of sin --- contravening the Prophetic directive - Ali and Fatima's marriage: the Prophet's own daughter married with nothing but a shield as mahr - The Qur'an promises enrichment: "If they be poor, Allah will enrich them of His bounty" (24:32) - The community must help the poor marry --- Umar ibn al-Khattab arranged his son's marriage and paid from the Public Treasury; zakat funds may be used for this purpose

\begin{discussionbox} 1. Why did the Prophet say "the best of marriages are the simplest ones"? \\ 2. What happens when weddings become too expensive? How does it affect young people? \\ 3. Ali married Fatima with nothing but a shield. What does this tell us about what matters most? \\ 4. How can our family support simpler, more Islamic celebrations? \end{discussionbox}

\begin{activitybox} Discuss the next family celebration (Eid, wedding, aqiqah). Make a concrete plan to simplify it following the Prophetic model. \end{activitybox}

\begin{duabox} "\arabicfont{اللَّهُمَّ} \arabicfont{اجْعَلْ} \arabicfont{خَيْرَ} \arabicfont{أَعْمَالِي} \arabicfont{خَوَاتِمَهَا}" --- "O Allah, make the best of my actions their endings." (variant du'a) \end{duabox}

\vspace{6pt}\noindent{\color{border}\rule{\textwidth}{0.5pt}}\vspace{6pt}

\subsubsection*{Module 3.4 --- Kindness to Parents}
\textbf{Week 16:} 

\begin{outcomebox} Family members understand that kindness to parents --- especially in old age --- is second only to worshipping Allah. \end{outcomebox}

\begin{quranbox} 17:23-24 --- "Your Lord has decreed that you worship none but Him and that you be kind to parents. If one or both of them attain old age with you, do not say a word of annoyance to them nor repulse them, but speak to them in gracious words and in mercy lower to them the wing of humility and say, 'My Lord, bestow Thy mercy on them, as they cherished me when I was little.'" \\ 46:15 --- "And We have enjoined on man kindness to his parents. His mother carries him in pain and she gives birth to him in pain, and (the period) of carrying him and weaning him is thirty months." \\ 31:14-15 --- "Be grateful to Me and to the parents; to Me is (the final) goal. But if they strive to compel thee to associate with Me that of which thou hast no knowledge, do not obey them; but keep company with them in this life in kind manner." \end{quranbox}

\begin{hadithbox} "Who is most deserving of my good companionship?" "Your mother," replied the Prophet (three times), then "Your father." (Bukhari, Muslim) \\ "Shall I not inform you about the three major sins? Associating partners with Allah, disobedience to parents, and telling lies and false testimony." (Bukhari, Muslim) \\ "Allah defers (the punishment of) all sins to the Day of Resurrection excepting disobedience to parents, for which Allah punishes the sinner in this life before his death." (al-Hakim) \\ "A man came to the Prophet and asked permission to go for jihad. The Prophet asked, 'Are your parents living?' 'Yes,' he replied. The Prophet then said, 'Then strive in their service.'" (Bukhari, Muslim, via Abdullah ibn Amr) \\ "Return to them and make them laugh as well as you made them weep." (Abu Dawud --- to the man who left his parents weeping) \end{hadithbox}

\paragraph*{Qaradawi's teaching (from *The Lawful and the Prohibited*, ch-03, pp.250-255)::}  - Disobedience to parents is a major sin, second only to shirk - The rights of the mother are stressed more --- because of suffering during pregnancy, childbirth, suckling, and rearing - Kindness especially in old age: "do not say 'uff' to them" --- the commentator said, "If a lesser thing than saying 'uff' were known, Allah would have prohibited it" - Parents' consent required even for jihad --- "go back to your parents and be a good companion to them" - Kindness even to non-Muslim parents: "keep company with them in this life in kind manner" (31:15) --- "such tolerant and beneficent teachings are not to be found in any other religion" - Causing parents to be insulted is a major sin: "Among the major sins is a man's cursing his parents"

\begin{discussionbox} 1. Why does Allah mention kindness to parents right after commanding worship of Him alone? \\ 2. The Prophet said "your mother" three times before "your father." Why? \\ 3. What does it mean to not even say "uff" to parents? How can we apply this with our parents/grandparents? \\ 4. What if parents ask us to do something wrong? How does the verse address this? \end{discussionbox}

\begin{activitybox} Call or visit grandparents this week. Each family member asks them one question about their life and listens attentively. Share what you learned at next halaqah. \end{activitybox}

\begin{duabox} "\arabicfont{رَبِّ} \arabicfont{ارْحَمْهُمَا} \arabicfont{كَمَا} \arabicfont{رَبَّيَانِي} \arabicfont{صَغِيرًا}" --- "My Lord, bestow mercy on them as they cherished me when I was little." (17:24) \end{duabox}

\vspace{6pt}\noindent{\color{border}\rule{\textwidth}{0.5pt}}\vspace{6pt}

\subsubsection*{Module 3.5 --- Children's Rights and Justice Among Children}
\textbf{Week 17:} 

\begin{outcomebox} Parents understand their responsibilities toward children, and the family understands the importance of fairness among siblings. \end{outcomebox}

\begin{quranbox} 17:31 --- "And do not kill your children out of fear of poverty; We shall provide for them and for you. Truly, the killing of them is a great sin." \\ 6:152 --- "And do not approach the property of an orphan except in the way that is best." \\ 4:10 --- "Indeed, those who devour the property of orphans unjustly are only consuming fire into their bellies." \end{quranbox}

\begin{hadithbox} "Each one of you is a caretaker (ra'i) and is responsible for those under his care." (Bukhari, Muslim) \\ "Wasting the sustenance of his dependants is sufficient sin for a man." (Ibn Hibban) \\ "Allah will ask every caretaker about what He entrusted him to take care of, whether he has done the care-taking well or otherwise, and (every) man will be asked concerning the people of his household." (Bukhari, Muslim) \\ "Do justice among your sons" --- repeated three times. (Bukhari, Muslim, via al-Nu'man ibn Bashir) \\ "Fear Allah and treat your children with equal justice." (Abu Dawud) \\ "Your children have the right of receiving equal treatment, as you have the right that they should honor you." (variant) \\ "I and the one who raises an orphan will be like these two in Paradise" --- and he pointed with his middle and index fingers. (Bukhari, Muslim) \end{hadithbox}

\paragraph*{Qaradawi's teaching (from *The Lawful and the Prohibited*, ch-03, pp.247-249)::}  - Children's rights: (1) right to life --- prohibition of killing; (2) right to a good name; (3) right to sustenance, education, and proper care; (4) right to lineage --- prohibition of attributing child to wrong father - "Wasting the sustenance of his dependants is sufficient sin for a man" --- providing is a religious obligation - Equal treatment of children obligatory, especially in gifts --- the story of Bashir ibn Sa'd: the Prophet refused to witness an unequal gift - "Do not ask me to be a witness to injustice. Your children have the right of receiving equal treatment" - Preferential treatment permitted only for valid reason (disability, need) --- per Imam Ahmad ibn Hanbal - Kindness to orphans: sponsoring an orphan (food, clothing, education, love) is a meritorious act --- "I and the one who raises an orphan will be like these two in Paradise" - Protecting orphan's wealth: devouring it unjustly is "consuming fire into their bellies" (4:10)

\begin{discussionbox} 1. What are the rights children have over their parents? What are the rights parents have over their children? \\ 2. The Prophet refused to witness an unequal gift. How can we make sure we are fair to all our children? \\ 3. What does "wasting the sustenance of dependants" mean? How can we be more responsible? \\ 4. How can our family support orphans in our community? \end{discussionbox}

\begin{activitybox} Parents audit their treatment of children: are gifts, attention, and time distributed fairly? Children share if they've ever felt treated unfairly. Discuss together how to fix it. \end{activitybox}

\begin{duabox} "\arabicfont{اللَّهُمَّ} \arabicfont{اجْعَلْ} \arabicfont{ذُرِّيَّتِي} \arabicfont{صَالِحِينَ}" --- "O Allah, make my offspring righteous." (variant du'a) \end{duabox}

\vspace{6pt}\noindent{\color{border}\rule{\textwidth}{0.5pt}}\vspace{6pt}

\subsubsection*{Module 3.6 --- Prophetic Pedagogy: Raising and Educating Children}
\textbf{Week 18:} 

\begin{outcomebox} Parents understand the Prophet's educational methods and commit to applying them in the home. \end{outcomebox}

\begin{quranbox} 66:6 --- "O you who believe, protect yourselves and your families from a Fire whose fuel is people and stones." \end{quranbox}

\begin{hadithbox} "A clever teacher is one who gives every human being his share of knowledge which suits him most and which is right for him to take." (Education and Economy in the Sunnah, ch-01) \\ "The people who receive the severest afflictions are the Prophets. Next to them is the selected group of Muslims and next to them the best of the remaining groups." (Bukhari, Muslim) \\ "Take full advantage of five before five: your youth before your old age, your health before your sickness, your wealth before your poverty, your life before your death, and your free time before your busy time." (al-Hakim, via Ibn Abbas) \end{hadithbox}

\paragraph*{Qaradawi's teaching (from *Education and Economy in the Sunnah*, ch-01)::}  - The Prophet established 19 educational principles that "anticipate many values attributed to modern pedagogy" - Gradual instruction: "Assuming a gradual attitude in teaching and making things easy rather than difficult" - Individual differences: "Taking individual differences and variations into consideration" --- the Prophet tailored counsel to each person's temperament and circumstances - Mercy and sympathy: "Being merciful and sympathetic towards the student" - Questions and dialogue: "Attracting the student's attention by asking questions and leading live conversation" - Stories and similitude: "Choosing the best modes (like similitude, giving examples, telling stories)" - Moderation: "Being moderate in the amount being taught and avoiding putting the student in a state of boredom" - Practical application: "Making good use of practical situations for the sake of education and guidance" - Rewarding good students; feeling pity toward mistaken ones (not anger) - "What is right for one person is not necessarily right for another, and what is suitable for a certain environment is not necessarily suitable for the other"

\begin{discussionbox} 1. How did the Prophet teach? Was he harsh or gentle? \\ 2. What does it mean to "take individual differences into consideration"? How are each of our children different? \\ 3. How can we use stories and questions instead of lectures in our family? \\ 4. What is one Prophetic teaching method we want to try this week? \end{discussionbox}

\begin{activitybox} Design a "family learning plan" --- for each child, identify: (1) their temperament, (2) their interests, (3) one Islamic lesson to teach them this week, and (4) the best method for teaching it based on their personality. \end{activitybox}

\begin{duabox} "\arabicfont{اللَّهُمَّ} \arabicfont{اجْعَلْنِي} \arabicfont{مُؤَدِّبًا} \arabicfont{حَسَنًا}" --- "O Allah, make me a good educator." (variant du'a) \end{duabox}

\vspace{6pt}\noindent{\color{border}\rule{\textwidth}{0.5pt}}\vspace{6pt}

\newpage
\section{PHASE 4: LEAVES --- LIVING ISLAM IN THE WORLD}

\begin{quotebox}\itshape "Your body has its rights, your wife has her rights, your guest has his rights." (Hadith). Faith and character and family all bear fruit in how we live daily --- with joy, with cleanliness, with purpose, and with healthy boundaries. This phase covers the practical application of Islam in everyday family life. Seven modules.\end{quotebox}

\subsubsection*{Module 4.1 --- Sabr: The Three Dimensions of Patience}
\textbf{Week 19:} 

\begin{outcomebox} Family members understand the three dimensions of sabr and identify which dimension is most needed in their life now. \end{outcomebox}

\begin{quranbox} 2:153 --- "Indeed, Allah is with the patient." \\ 91:9-10 --- "Successful is the one who purifies it, and failed is the one who corrupts it." \\ 29:69 --- "Those who strive for Us --- We will surely guide them to Our ways." \end{quranbox}

\begin{hadithbox} (Sabr appears 90+ times in the Qur'an --- the most frequently commanded virtue) \end{hadithbox}

\paragraph*{Qaradawi's teaching (from *concept-sabr*; *Ethics in Islam* ch-03, ch-04)::}  - Three dimensions: 1. \textbf{Sabr fi al-Ta'ah} --- patience IN obedience: persisting in prayer, fasting, zakat, and good deeds even when difficult 2. \textbf{Sabr 'an al-Ma'siyah} --- patience AGAINST sin: steadfastly refusing temptation and resisting the pull of wrongdoing 3. \textbf{Sabr 'ala al-Musibah} --- patience IN affliction: enduring trials, loss, and hardship without despair --- paired with hope (raja') in Allah - Sabr is paired with hope: the patient person does not just endure but trusts that relief is coming - "The hearts get bored just as the bodies do" --- patience includes patience with the monotony of daily life

\begin{discussionbox} 1. Which of the three types of sabr is hardest for you? (patience in worship, patience against temptation, patience in hardship) \\ 2. Which one do we as a family need most right now? \\ 3. What does "Allah is with the patient" mean? How does knowing Allah is with you make patience easier? \\ 4. Can you share a time when being patient led to something good? \end{discussionbox}

\begin{activitybox} "Sabr tracker" --- each family member picks one area where they need sabr this week. Track daily. Share results at next halaqah. \end{activitybox}

\begin{duabox} "\arabicfont{اللَّهُمَّ} \arabicfont{إِنِّي} \arabicfont{أَسْأَلُكَ} \arabicfont{الصَّبْرَ} \arabicfont{وَالْيَقِينَ}" --- "O Allah, I ask You for patience and certainty." (variant du'a) \end{duabox}

\vspace{6pt}\noindent{\color{border}\rule{\textwidth}{0.5pt}}\vspace{6pt}

\subsubsection*{Module 4.2 --- Wasatiyyah: The Middle Way in Family Life}
\textbf{Week 20:} 

\begin{outcomebox} Family members understand that wasatiyyah is not compromise but the straight path between two extremes of error. \end{outcomebox}

\begin{quranbox} 2:143 --- "And thus We have made you a middle nation, that you will be witnesses over the people, and the Messenger will be a witness over you." \\ 25:67 --- "Who when they spend, are not extravagant nor niggardly, but find a middle way between." \end{quranbox}

\begin{hadithbox} "Your body has its rights, your wife has her rights, your guest has his rights. Give each their rights." (Bukhari 1975, Muslim 1159) \\ The Prophet corrected three men who vowed extreme asceticism: fasting every day, praying all night, abstaining from marriage. \end{hadithbox}

\paragraph*{Qaradawi's teaching (from *concept-wasatiyyah*; *comparison-wasatiyyah-moderation*)::}  - Wasatiyyah is Qaradawi's master framework --- it structures his thought across fiqh, ethics, lifestyle, and spirituality - Wasatiyyah is NOT compromise between truth and falsehood --- it is the straight path between two extremes of ERROR - Three applications: 1. \textbf{In worship:} Not excessive asceticism, not mere ritualism --- balance: worship + family + rest + social duties 2. \textbf{In fiqh:} Not blind taqlid, not reckless innovation --- renewed ijtihad grounded in authentic sources 3. \textbf{In ethics:} Not extremism, not secularism --- Islam regulates freedom and keeps desire in check - The Prophet rejected: fasting every day, praying all night, abstaining from marriage - Moderation in emotions: "Love moderately, for you may hate those you love one day. And hate moderately, for you may love those you hate one day."

\begin{discussionbox} 1. What does "the middle way" mean for our family? Where do we tend to go to extremes? \\ 2. The Prophet said "your body has its rights." How do we balance worship with rest and family? \\ 3. Wasatiyyah is not compromise between right and wrong. What's the difference? \\ 4. What is one area where our family needs more balance? \end{discussionbox}

\begin{activitybox} "Family balance audit" --- list all weekly activities. Are any areas extreme (too much screen time, too much work, too much socializing, too much isolation)? Make one adjustment. \end{activitybox}

\begin{duabox} "\arabicfont{اللَّهُمَّ} \arabicfont{اجْعَلْنَا} \arabicfont{مِنْ} \arabicfont{أُمَّةِ} \arabicfont{الْوَسَطِ}" --- "O Allah, make us among the middle nation." (variant du'a from Q 2:143) \end{duabox}

\vspace{6pt}\noindent{\color{border}\rule{\textwidth}{0.5pt}}\vspace{6pt}

\subsubsection*{Module 4.3 --- Joy, Humor, and Celebration in Family Life}
\textbf{Week 21:} 

\begin{outcomebox} Family members understand that joy and laughter are part of faith, following the Prophet's example. \end{outcomebox}

\begin{quranbox} 90:4 --- "Indeed, We have created mankind in struggle." \\ 7:31-32 --- "O children of Adam, take your adornment at every masjid... Who has forbidden the adornment of Allah which He has produced for His servants?" \end{quranbox}

\begin{hadithbox} "Laughter distinguishes human beings from other species." (Qaradawi's framing, Diversion and Arts ch-07) \\ The Prophet jested "without resorting to lies" --- he shared companions' laughter, joys, and jests (Zayd ibn Thabit's testimony) \\ The Prophet bent down so his grandsons (Hasan and Hussein) would ride on his back. When a companion said "You ride the back of the best of men," he said: "They are the best of riders!" (Bukhari) \\ The Prophet jested with an old woman about Paradise: "No old woman enters Paradise!" --- then explained she would enter as a young woman (56:35-36) \\ Aisha's food-fight with Sawda: Aisha smeared Sawda's face with food; the Prophet lowered his knees so Sawda could get even, and "the Prophet stood there laughing" (variant hadith) \\ "O Abu Bakr! Let them sing for it is feast day." (Bukhari --- the Prophet defending singing at Eid) \\ "Haven't you got any amusement for the marriage ceremony? The Ansar like amusement." (Ibn Majah) \\ "The difference between what is lawful and what is unlawful is the duff and the voice at weddings." (Bukhari) \end{hadithbox}

\paragraph*{Qaradawi's teaching (from *Diversion and Arts in Islam*, ch-07; *concept-art*)::}  - "Life is a journey of hardships shrouded in toil and trouble" --- people need relief - "Islam is the religion of instinct; it cannot possibly repress the human urge to laugh and have a good time" - The Prophet "used to jest without resorting to lies" --- he "lived a normal life with his companions. He shared their laughters, joys and jests. He also shared their pain, grief and trials." - Companions described him as "one of the most humorous men" - "At home, he used to jest and have fun with his wives. He even listened to their short stories." - The Prophet liked joyous occasions: Eids and weddings --- he encouraged celebration with song and amusement - Omar ibn al-Khattab, known for sternness, also jested with his slave girl

\begin{discussionbox} 1. Was the Prophet serious all the time? What do the hadiths tell us? \\ 2. How can laughter and joy be part of faith? \\ 3. The Prophet played with his grandchildren. How do we play with ours? \\ 4. How can we celebrate Eid and family milestones in a way that is joyful AND Islamic? \end{discussionbox}

\begin{activitybox} Family game night --- play a game, tell jokes (clean ones!), share funny family memories. End with du'a of gratitude. \end{activitybox}

\begin{duabox} "\arabicfont{اللَّهُمَّ} \arabicfont{اجْعَلْنِي} \arabicfont{مِنَ} \arabicfont{الْمُتَوَكِّلِينَ} \arabicfont{عَلَيْكَ} \arabicfont{الْمُسْتَبْشِرِينَ} \arabicfont{بِرَحْمَتِكَ}" --- "O Allah, make me among those who trust in You and rejoice in Your mercy." (variant du'a) \end{duabox}

\vspace{6pt}\noindent{\color{border}\rule{\textwidth}{0.5pt}}\vspace{6pt}

\subsubsection*{Module 4.4 --- Beauty and Aesthetics in Islam}
\textbf{Week 22:} 

\begin{outcomebox} Family members understand that beauty is a divine attribute and that appreciating beauty is part of faith. \end{outcomebox}

\begin{quranbox} 16:5-6 --- "God created cattle for you (you find) in them warmth, benefit and from them you eat. And in them there is beauty for you when you bring them home and when you take them to pasture." \\ 16:8 --- "God created horses, mules and donkeys for you to ride and for ornament." \\ 16:14 --- "God is the One Who has subjected to you the sea, so that you eat from it fresh flesh, and you extract from it ornaments which you wear." \\ 6:141 --- "Look at their fruits when they bear fruit, and their ripening. Surely in that there are signs for people who believe." \\ 7:31-32 --- "O children of Adam, take your adornment at every masjid... Who has forbidden the adornment of Allah which He has produced for His servants?" \end{quranbox}

\begin{hadithbox} "Allah is Beautiful and He loves beauty." (Muslim, via Ibn Mas'ud) \\ When a companion said "All people like to wear attractive clothes and shoes," the Prophet replied: "Allah is Beautiful and He loves beauty, but arrogance is different. It means to deny the truth although one is certain that it is so, and to be thankless for the people who were helpful to one." (Muslim) \\ "Adorn the Qur'an by reciting it with your nice voices." (al-Bukhari and others) \\ "Cleanliness is half of one's faith." (Muslim) \end{hadithbox}

\paragraph*{Qaradawi's teaching (from *Diversion and Arts in Islam*, ch-01, ch-02; *concept-art*; *concept-taqwa*)::}  - Islam does not reject art or beauty --- it channels them - The Qur'an consistently pairs aesthetic appreciation with utilitarian benefit: cattle provide warmth AND beauty; horses for riding AND ornament; the sea for food AND jewelry; fruits for nutrition AND visual delight - "The Qur'an orientates man towards the assimilation of aesthetic and utilitarian aspects in different fields of life" - "Allah is Beautiful and He loves beauty" --- appreciation of beauty is part of faith, not opposed to it - The Qur'an itself is "both an aesthetic and intellectual miracle" --- "It fulfils man's spiritual, intellectual and emotional needs" - The Prophet corrected confusion between love of beauty and arrogance: beauty ≠ pride - Art serves human needs: "Sports keep the body fit, worship purifies the soul, knowledge develops the intellect, and art is a source of emotional fulfilment" - Seven key principles governing art: (1) default is permissibility; (2) content over form; (3) middle course; (4) art serves human needs; (5) art within Islamic framework; (6) beauty is a divine attribute; (7) prevention of idolatry

\begin{discussionbox} 1. What does "Allah is Beautiful and He loves beauty" mean for our daily life? \\ 2. The Qur'an pairs beauty AND benefit --- cattle, horses, the sea, fruits. Can you think of something that is both beautiful and useful? \\ 3. How can we make our home more beautiful in a way that pleases Allah? \\ 4. What is the difference between loving beauty and being arrogant? \end{discussionbox}

\begin{activitybox} Family beauty walk --- go outside together and find 5 beautiful things in nature. Each person explains why they find it beautiful. Connect to "Allah is Beautiful and He loves beauty." \end{activitybox}

\begin{duabox} "\arabicfont{اللَّهُمَّ} \arabicfont{أَنْتَ} \arabicfont{الْجَمِيلُ} \arabicfont{وَتُحِبُّ} \arabicfont{الْجَمَالَ}" --- "O Allah, You are Beautiful and You love beauty." (Muslim, via Ibn Mas'ud) \end{duabox}

\vspace{6pt}\noindent{\color{border}\rule{\textwidth}{0.5pt}}\vspace{6pt}

\subsubsection*{Module 4.5 --- Cleanliness: Half of Faith}
\textbf{Week 23:} 

\begin{outcomebox} Family members understand that cleanliness is a form of worship and develop concrete hygiene habits. \end{outcomebox}

\begin{quranbox} 2:222 --- "Indeed, Allah loves those who are constantly repentant and loves those who purify themselves." \\ 74:4 --- "And your clothing purify." \\ 7:31 --- "O children of Adam, take your adornment at every masjid." \end{quranbox}

\begin{hadithbox} "Cleanliness is half of one's faith." (Muslim) \\ "Allah is clean, and He loves cleanliness. Clean your courtyards, and do not imitate the Jews." (Tirmidhi, via Sa'd ibn Abi Waqqas) \\ "The Miswak (tooth-stick) cleanses and purifies the mouth and pleases Allah." (Ahmed, al-Nassa'i, via Aisha) \\ "If you have hair, honor it." (Abu Dawud, via Abu Hurayrah) \\ "Removing harm from the road is a form of charity." (Bukhari, Muslim, via Abu Hurayrah) \\ "Allah does not look at your appearances and your wealth, but He looks at your hearts and your deeds." (Muslim) \end{hadithbox}

\paragraph*{Qaradawi's teaching (from *Ethics in Islam*, section-04-02, pp.338-340)::}  - Cleanliness plays a significant role in faith --- "faith is cleansing oneself from evil and developing noble qualities" - Cleanliness linked to worship: one cannot pray unless physically clean; "almost all books on Islamic jurisprudence start with a chapter on cleanliness" - The first verses revealed include: "And your clothing purify" (74:4) - Personal hygiene practices: brushing teeth (miswak), washing hands before and after meals, removing excess hair, trimming nails - Islam encourages adornment with fine clothes and perfumes, especially at Hajj and religious festivities - Keeping houses, mosques, and streets clean --- "Allah is clean, and He loves cleanliness" - Modesty, humility, and cleanliness are interconnected: inner purity expressed through outer cleanliness

\begin{discussionbox} 1. The Prophet said "cleanliness is half of faith." How is cleaning yourself an act of worship? \\ 2. What are the daily cleanliness habits the Sunnah teaches us? \\ 3. How can we keep our home and neighborhood cleaner as a family? \\ 4. What does inner cleanliness (purifying the heart) look like? \end{discussionbox}

\begin{activitybox} Family cleaning day --- everyone cleans one area of the home together. Frame it as worship, not a chore. Make du'a together afterward. \end{activitybox}

\begin{duabox} "\arabicfont{اللَّهُمَّ} \arabicfont{اجْعَلْنِي} \arabicfont{مِنَ} \arabicfont{التَّوَّابِينَ} \arabicfont{وَاجْعَلْنِي} \arabicfont{مِنَ} \arabicfont{الْمُتَطَهِّرِينَ}" --- "O Allah, make me among those who repent and make me among those who purify themselves." (variant from Q 2:222) \end{duabox}

\vspace{6pt}\noindent{\color{border}\rule{\textwidth}{0.5pt}}\vspace{6pt}

\subsubsection*{Module 4.6 --- Wise Use of Time}
\textbf{Week 24:} 

\begin{outcomebox} Family members understand that time is life and that we are accountable for how we spend it. \end{outcomebox}

\begin{quranbox} 103:1-3 --- "By time, indeed, mankind is in loss, except for those who believe and do righteous deeds and join together in the truth and join together in patience." \\ 63:10-11 --- "Spend [in the way of Allah] from what We have provided you before death approaches one of you, and he says, 'My Lord, if only You would delay me for a brief term so I would give to charity and be among the righteous.' But never will Allah delay a soul when its time has come." \\ 16:12 --- "He has subjected for you the night and day." \end{quranbox}

\begin{hadithbox} "One's feet will not move on the Day of Resurrection before he is asked about his life, how he spent it; his youth, how he wore it out; his wealth, how he earned it and how he spent it; and his knowledge, how he used it." (Tirmidhi, via Ibn Mas'ud) \\ "Take full advantage of five before five: your youth before your old age, your health before your sickness, your wealth before your poverty, your life before your death, and your free time before your busy time." (al-Hakim, via Ibn Abbas) \\ "Whoever begins the day feeling secure with his family and healthy with enough provisions for the day is as though he owns the whole world." (Tirmidhi, via Ubaydullah ibn Muhsin) \end{hadithbox}

\paragraph*{Qaradawi's teaching (from *Ethics in Islam*, section-04-02, pp.344-345)::}  - The Qur'an draws attention to the value of time through divine oaths: "By the night" (92:1), "By the dawn" (89:1), "By the morning brightness" (93:1), "By time" (103:1) - Qur'an shows people facing death wishing they had used time better (63:10-11; 35:37) - The Prophet: "One's feet will not move on the Day of Resurrection before he is asked about his life, how he spent it" - "Take full advantage of five before five" --- the Prophetic framework for seizing time - Hassan al-Banna: "Our life is nothing but the time we spend between our birth and our death" - Imam al-Hasan al-Basri: "You are nothing but a limited number of days. With every day that passes, part of you is gone." - Traditional saying: "Every day, as dawn breaks, dawn calls on people: 'O children of Adam, I am a new creation and a witness to what you do. Take what you can from me, because once I am gone, I will never be back.'" - Muslims should focus on acts of worship at special times (dawn, Friday, Ramadan) when rewards are multiplied

\begin{discussionbox} 1. The Qur'an swears "By Time." What does that tell us about how important time is? \\ 2. What will we be asked about on the Day of Resurrection? How does that make you think about your time? \\ 3. "Take advantage of five before five." Which one is most relevant to you right now? \\ 4. What is one way our family wastes time? How can we change it? \end{discussionbox}

\begin{activitybox} Family time audit --- track how each person spends one full day. Review together. Identify: what time is well spent? What is wasted? Make one family change. \end{activitybox}

\begin{duabox} "\arabicfont{اللَّهُمَّ} \arabicfont{بَارِكْ} \arabicfont{لَنَا} \arabicfont{فِي} \arabicfont{أَوْقَاتِنَا}" --- "O Allah, bless us in our time." (variant du'a) \end{duabox}

\vspace{6pt}\noindent{\color{border}\rule{\textwidth}{0.5pt}}\vspace{6pt}

\subsubsection*{Module 4.7 --- Protecting the Family: Boundaries and Halal/Haram}
\textbf{Week 25:} 

\begin{outcomebox} Family members understand that Islam sets boundaries to protect the family and that these boundaries come from Allah's wisdom. \end{outcomebox}

\begin{quranbox} 17:32 --- "And do not come near zina; indeed, it is an abomination and an evil way." \\ 24:30-31 --- "Tell the believing men that they should lower their gazes and guard their private parts... And tell the believing women that they should lower their gazes and guard their private parts and not display their adornment except that which appears thereof." \end{quranbox}

\begin{hadithbox} "Whoever believes in Allah and the Last Day must never be in privacy with a woman without there being a mahram with her, for otherwise Satan will be the third person." (Ahmad, via Amir ibn Raby'ah) \\ "The in-law is death." (Bukhari, Muslim --- the Prophet's warning about male in-laws being alone with one's wife) \\ "Do not give a second look. While you are not to blame for the first, you have no right to the second." (Ahmed, via Ali ibn Abi Talib) \\ "The difference between what is lawful and what is unlawful is the duff and the voice at weddings." (Bukhari) \end{hadithbox}

\paragraph*{Qaradawi's teaching (from *The Lawful and the Prohibited*, ch-03, pp.168-171; *Ethics in Islam* section-04-02, pp.337-338)::}  - Three approaches to desire: (1) free sex --- destroys family and society; (2) suppression --- contradicts nature and Allah's plan; (3) regulation --- Islam's middle path - Islam doesn't just prohibit zina --- it blocks every avenue leading to it: "Do not come near zina" (17:32) - Khalwah prohibition: a man and non-mahram woman must not be alone --- "Satan will be the third" - Special warning about in-laws: "The in-law is death" --- because relatives have easier access and people are negligent - Lowering the gaze: BOTH men and women commanded --- "the eye is the key to the feelings, and the look is a messenger of desire" - Chastity (iffah): not celibacy, but lawful enjoyment within marriage; "safeguard their private parts --- except from their wives" (23:5-7) - "Do not give a second look" --- the first is accidental, the second is a choice

\begin{discussionbox} 1. Why does Islam say "do not come NEAR zina" instead of just "do not commit zina"? What's the difference? \\ 2. What does "lowering the gaze" mean in the age of phones and screens? \\ 3. How can we protect our family from harmful influences online and in social situations? \\ 4. What is the difference between a boundary and a restriction? Are boundaries protection or limitation? \end{discussionbox}

\begin{activitybox} Family media agreement --- discuss and agree on 3 family rules for phones, internet, and social situations. Write them down. Everyone signs. \end{activitybox}

\begin{duabox} "\arabicfont{اللَّهُمَّ} \arabicfont{اسْتُرْ} \arabicfont{عَوْرَاتِنَا} \arabicfont{وَأَحْفَظْنَا} \arabicfont{مِنْ} \arabicfont{بَيْنِ} \arabicfont{أَيْدِينَا} \arabicfont{وَمِنْ} \arabicfont{خَلْفِنَا}" --- "O Allah, conceal our faults and protect us from front and behind." (variant from Q 17:80) \end{duabox}

\vspace{6pt}\noindent{\color{border}\rule{\textwidth}{0.5pt}}\vspace{6pt}

\newpage
\section{PHASE 5: FRUIT --- TRUST, HOPE, AND ETERNAL VISION}

\begin{quotebox}\itshape "Whoever fears Allah --- He will make for him a way out." (Qur'an 65:2). The fruit of faith, character, family, and daily practice is a family that trusts Allah completely, hopes in His mercy, and lives with eternal vision. Five modules.\end{quotebox}

\subsubsection*{Module 5.1 --- Justice in the Family and Beyond}
\textbf{Week 26:} 

\begin{outcomebox} Family members understand that justice is a divine command and that it begins at home. \end{outcomebox}

\begin{quranbox} 16:90 --- "Indeed, Allah orders justice and good conduct and giving to relatives and forbids immorality and bad conduct and oppression." \\ 4:135 --- "Stand out firmly for justice, as witnesses to Allah, even if against yourselves or parents or near relatives." \\ 5:8 --- "Be just; that is nearer to righteousness." \\ 4:58 --- "Indeed, Allah commands you to render trusts to whom they are due, and when you judge between people, to judge with justice." \end{quranbox}

\begin{hadithbox} "The most beloved of people to Allah on the Day of Judgment is the just leader, and the most hated of people to Allah on the Day of Judgment is the unjust leader." (Tirmidhi) \end{hadithbox}

\paragraph*{Qaradawi's teaching (from *concept-justice*; *Ethics in Islam* section-01-03, ch-02; *Faith and Life* ch-01)::}  - Justice is not merely a virtue --- it is a binding divine command: "Allah orders justice" (16:90) - Three dimensions: (1) justice toward Allah --- worship and obedience; (2) justice toward others --- rights, transactions, governance; (3) justice toward oneself --- moderation, self-discipline - "Stand out firmly for justice, as witnesses to Allah, even if against yourselves or parents or near relatives" (4:135) --- justice even when it costs you personally - Justice requires an Islamic framework --- penalties cannot be applied in isolation from social welfare and economic justice - "In a true Muslim society, no one steals out of poverty" --- justice is systemic, not just individual - Economic justice: zakat as redistribution; "And in their wealth there is a right of the needy" (51:19) - Justice rooted in human dignity: every person honored by Allah (17:70) --- racism, classism, and discrimination are violations

\begin{discussionbox} 1. What does justice look like in our family? Between spouses? Between siblings? Between parents and children? \\ 2. "Stand out firmly for justice even if against yourselves." What does that mean practically? \\ 3. How can we be more just in how we treat extended family, neighbors, and community? \\ 4. Is there any injustice we need to correct in our family? \end{discussionbox}

\begin{activitybox} Family justice circle --- each member shares one time they felt treated unfairly in the family. Parents listen without defending. Discuss how to make it right. \end{activitybox}

\begin{duabox} "\arabicfont{اللَّهُمَّ} \arabicfont{أَعِنِّي} \arabicfont{عَلَى} \arabicfont{الْعَدْلِ} \arabicfont{حِينَ} \arabicfont{أَرْضَى} \arabicfont{وَحِينَ} \arabicfont{أَغْضَبُ}" --- "O Allah, help me to be just whether I am pleased or angry." (variant du'a) \end{duabox}

\vspace{6pt}\noindent{\color{border}\rule{\textwidth}{0.5pt}}\vspace{6pt}

\subsubsection*{Module 5.2 --- Responsibility and Charity}
\textbf{Week 27:} 

\begin{outcomebox} Family members understand that every person is accountable before Allah and that charity is a family obligation. \end{outcomebox}

\begin{quranbox} 6:164 --- "Every soul earns not [blame] except against itself, and no bearer of burdens will bear the burden of another." \\ 33:72 --- "Indeed, We offered the trust to the heavens and the earth and the mountains, and they declined to bear it and feared it; but man bore it." \\ 51:19 --- "And in their wealth there is a right of the needy, he who asks, and he who is deprived." \\ 2:215 --- "Whatever you spend of good is for parents and relatives and orphans and the needy and the traveler." \end{quranbox}

\begin{hadithbox} "Each one of you is a caretaker (ra'i) and is responsible for those under his care." (Bukhari, Muslim) \\ "A believer to another believer is like a building whose different parts enforce each other." (Bukhari, Muslim, via Abu Musa al-Ash'ari) \\ "Charity is not permissible to be given to a rich person, or for one who is strong and healthy." (Ahmed, Abu Dawud, via Abdullah ibn Amr) \\ "Fire is more appropriate for a body that grew on ill-gotten money." (Ahmed, Tirmidhi, via Jabir ibn Abdullah) \end{hadithbox}

\paragraph*{Qaradawi's teaching (from *concept-responsibility*; *concept-zakat*; *Ethics in Islam* section-04-02; *Economic Security* ch-04, ch-06)::}  - Individual responsibility is non-transferable: "No person can carry the sin of another" - Communal responsibility complements individual duty: the wealthy are responsible for the poor; scholars for guidance; rulers for justice - The trust of responsibility: "the heavens and the earth and the mountains refused the trust, and man bore it" (33:72) - Prophetic model integrates "freedom and responsibility, individualism and collectivism" - Charity: zakat as obligatory redistribution; sadaqah as voluntary; spending on family as charity - "Charity is not permissible for a rich person or one who is strong and healthy" --- charity goes to those who truly need it - Giving to relatives first: "Whatever you spend of good is for parents and relatives and orphans and the needy" - Ill-gotten money: "Fire is more appropriate for a body that grew on ill-gotten money"

\begin{discussionbox} 1. What does it mean that "no bearer of burdens will bear the burden of another"? How does this apply to our family? \\ 2. "Each one of you is a caretaker." Who is each person in our family a caretaker of? \\ 3. How can our family give charity together? Who in our community needs help? \\ 4. What does it mean that the mountains "refused the trust" but humans accepted it? \end{discussionbox}

\begin{activitybox} Family sadaqah project --- choose one act of charity to do together this month (cook a meal for someone, donate clothes, visit a sick person, give to a relative in need). \end{activitybox}

\begin{duabox} "\arabicfont{اللَّهُمَّ} \arabicfont{اجْعَلْنِي} \arabicfont{مُنْفِقًا} \arabicfont{فِي} \arabicfont{سَبِيلِكَ}" --- "O Allah, make me one who spends in Your way." (variant du'a) \end{duabox}

\vspace{6pt}\noindent{\color{border}\rule{\textwidth}{0.5pt}}\vspace{6pt}

\subsubsection*{Module 5.3 --- Inner Struggle: Jihad al-Nafs}
\textbf{Week 28:} 

\begin{outcomebox} Family members understand that the greatest jihad is the struggle against one's own desires, and that this is a lifelong, daily effort. \end{outcomebox}

\begin{quranbox} 91:7-10 --- "And by the soul and He who proportioned it --- And inspired it [with discernment of] its wickedness and its righteousness. He has succeeded who purifies it. And he has failed who instills it [with corruption]." \\ 79:40-41 --- "He who fears the position of his Lord and prevents the soul from [unlawful] inclination --- Then indeed, Paradise will be [his] refuge." \\ 29:69 --- "Those who strive for Us --- We will surely guide them to Our ways. And indeed, Allah is with the doers of good." \\ 38:26 --- "Oh David, indeed We have made you a successor upon the earth, so judge between people in truth and do not follow [your own] desire, as it will lead you astray from the way of Allah." \end{quranbox}

\begin{hadithbox} (The inner struggle is described in the tradition as the "greater jihad") \end{hadithbox}

\paragraph*{Qaradawi's teaching (from *Ethics in Islam*, section-04-02, pp.342-343)::}  - "Human nature is just as capable of doing evil as it is of doing good. It is always easier to slip to the evil side than to elevate oneself to purity." - "The hardest obstacles on one's journey to self-purification is the inner self with its desires and tendency to give in to temptation" - Qur'an warns: "Who is more astray than one who follows his desire without guidance from Allah?" (28:50) - "If We had willed, we could have elevated him thereby, but he adhered instead to the earth and followed his own desire. His example is like that of the dog" (7:176) - To succeed: "One has to be vigilant, but patient with oneself" - "He who fears the position of his Lord and prevents the soul from [unlawful] inclination --- Paradise will be his refuge" (79:40-41) - "This inner struggle is in fact more difficult than fighting enemies. This is why it is often described in the tradition as the 'greater jihad'" - "Those who strive for Us --- We will surely guide them to Our ways" (29:69) --- the promise of guidance for those who struggle

\begin{discussionbox} 1. What is the "greater jihad"? How is it harder than fighting? \\ 2. What is one desire or habit that you struggle with? (share age-appropriately) \\ 3. The Qur'an says "it is always easier to slip to the evil side." What helps us resist? \\ 4. How can our family support each other in the inner struggle? \end{discussionbox}

\begin{activitybox} "One habit challenge" --- each family member identifies one bad habit they want to overcome this month. The family supports each person. Check in weekly. \end{activitybox}

\begin{duabox} "\arabicfont{اللَّهُمَّ} \arabicfont{أَعِنِّي} \arabicfont{عَلَى} \arabicfont{ذَكْرِكَ} \arabicfont{وَشُكْرِكَ} \arabicfont{وَحُسْنِ} \arabicfont{عِبَادَتِكَ}" --- "O Allah, help me to remember You, thank You, and worship You in the best way." (Abu Dawud, via Mu'adh ibn Jabal) \end{duabox}

\vspace{6pt}\noindent{\color{border}\rule{\textwidth}{0.5pt}}\vspace{6pt}

\subsubsection*{Module 5.4 --- Dawah Within the Family}
\textbf{Week 29:} 

\begin{outcomebox} Family members understand that dawah begins at home and that the best dawah is living example. \end{outcomebox}

\begin{quranbox} 16:125 --- "Invite to the way of your Lord with wisdom and good instruction, and argue with them in a way that is best." \\ 33:36 --- "It is not for a believing man or a believing woman, when Allah and His Messenger have decided a matter, that they should [thereafter] have any choice about their affair." \end{quranbox}

\begin{hadithbox} "A believer to another believer is like a building whose different parts enforce each other." (Bukhari, Muslim, via Abu Musa al-Ash'ari) \end{hadithbox}

\paragraph*{Qaradawi's teaching (from *concept-dawah*; *concept-wisdom*; *Approaching the Sunnah* ch-01)::}  - Dawah = inviting others to Islam --- including family members who may need renewal - Dawah must renew understanding, not alter the religion: "the caller invites people back to the original religion as practiced by the Prophet and Companions --- not to a 'new edition' of Islam" - Renewal means "revival of ijtihad in it, and recourse to its original wellsprings, liberation from inflexibility and imitativeness" - Dawah pursues the middle way: the Prophetic model rejects both extremism and laxity - Wisdom (hikmah) is the primary tool: "Invite to the way of your Lord with wisdom and good instruction" (16:125) - Dawah includes defending Islam against misrepresentation - The Prophetic pattern is "a comprehensive pattern" that integrates "freedom and responsibility"

\begin{discussionbox} 1. What is dawah? Is it only for scholars, or is it for everyone? \\ 2. How can we do dawah within our own family --- helping each other grow in faith? \\ 3. What does "invite with wisdom" mean? How is that different from lecturing or nagging? \\ 4. What is the best form of dawah? (Answer: living example --- being a model of what you teach) \end{discussionbox}

\begin{activitybox} "Family dawah plan" --- each family member identifies one person outside the family they can help toward Islam through good character and kind words (not preaching). Commit to one action this month. \end{activitybox}

\begin{duabox} "\arabicfont{اللَّهُمَّ} \arabicfont{اجْعَلْنِي} \arabicfont{دَاعِيًا} \arabicfont{إِلَى} \arabicfont{سَبِيلِكَ} \arabicfont{بِحِكْمَةٍ}" --- "O Allah, make me a caller to Your way with wisdom." (variant from Q 16:125) \end{duabox}

\vspace{6pt}\noindent{\color{border}\rule{\textwidth}{0.5pt}}\vspace{6pt}

\subsubsection*{Module 5.5 --- Trust in Allah's Plan: The Family That Surrenders}
\textbf{Week 30:} 

\begin{outcomebox} The family consolidates the entire course --- reviewing the journey from faith to character to family to daily life --- and commits to living as a family that trusts Allah completely. \end{outcomebox}

\begin{quranbox} 65:2-3 --- "Whoever fears Allah --- He will make for him a way out. And will provide for him from where he does not expect." \\ 3:159 --- "So by mercy from Allah, [O Muhammad], you were lenient with them. And if you had been rude [in speech] and harsh in heart, they would have disbanded from about you. So pardon them and ask forgiveness for them and consult them in the matter." \\ 2:156-157 --- "Who, when disaster strikes them, say, 'Indeed we belong to Allah, and indeed to Him we will return.' Those are the ones upon whom are blessings from their Lord and mercy." \end{quranbox}

\begin{hadithbox} "Tie it and rely on Allah." (Tirmidhi, via Anas ibn Malik) \\ "Oh Abu Bakr! What do you think of two men in the company of Allah? Have no fear, for Allah is with us." (Bukhari, via Abu Bakr) \\ "Oh Allah, I beseech You for guidance, piety, chastity, and self-sufficiency." (Muslim, via Tirmidhi --- the Prophet's comprehensive du'a) \end{hadithbox}

\paragraph*{Qaradawi's teaching (synthesis across all books)::}  - This module synthesizes the entire course: - Faith (Phase 1) gives us the foundation --- we are Allah's honored vicegerents - Character (Phase 2) gives us the trunk --- good morals, modesty, humility, sincerity - Family (Phase 3) gives us the branches --- marriage, parent-child, Prophetic pedagogy - Daily life (Phase 4) gives us the leaves --- patience, moderation, joy, cleanliness, time, boundaries - Trust (Phase 5) gives us the fruit --- a family that surrenders to Allah completely - The Qur'an's promise: "Whoever fears Allah --- He will make for him a way out. And will provide for him from where he does not expect" (65:2-3) --- this is the ultimate fruit of taqwa - "Indeed we belong to Allah, and indeed to Him we will return" (2:156) --- the family's creed in both joy and trial - Tawakkul is the culmination: tie your camel, prepare everything, AND trust Allah completely --- both/and, not either/or - The Prophet's du'a brings everything together: guidance, piety (taqwa), chastity ('afaf), and self-sufficiency (ghina)

\begin{discussionbox} 1. What did we learn over these 30 weeks? What changed in our family? \\ 2. Which module was most impactful for each person? \\ 3. The Qur'an promises "whoever fears Allah --- He will make for him a way out." How have we seen this in our family's life? \\ 4. What will we commit to continuing after this course ends? \end{discussionbox}

\begin{activitybox} Family covenant --- write a family mission statement based on the 5 phases: - "Our family is rooted in FAITH, grows through CHARACTER, branches as FAMILY, lives in the WORLD with balance and joy, and bears the FRUIT of trust in Allah." \end{activitybox}

Sign it, frame it, hang it. Plan the next halaqah cycle.

\begin{duabox} "\arabicfont{رَبَّنَا} \arabicfont{هَبْ} \arabicfont{لَنَا} \arabicfont{مِنْ} \arabicfont{أَزْوَاجِنَا} \arabicfont{وَذُرِّيَّاتِنَا} \arabicfont{قُرَّةَ} \arabicfont{أَعْيُنٍ} \arabicfont{وَاجْعَلْنَا} \arabicfont{لِلْمُتَّقِينَ} \arabicfont{إِمَامًا}" --- "Our Lord, grant us from among our spouses and offspring comfort to our eyes, and make us a leader for the righteous." (25:74) \end{duabox}

\vspace{6pt}\noindent{\color{border}\rule{\textwidth}{0.5pt}}\vspace{6pt}

\subsection*{Course Summary}

\begin{longtable}{|p{2.5cm}|p{3.5cm}|p{4.5cm}|p{5cm}|}
\hline
\textbf{\color{accent}Phase} & \textbf{\color{accent}Theme} & \textbf{\color{accent}Weeks} & \textbf{\color{accent}Modules} \\
\hline
1. Roots & Faith as the Foundation & 1–6 & Dignity, Happiness, Love, Hope, Taqwa, Tawakkul \\
\hline
2. Trunk & Character \& Purification & 7–12 & Akhlaq, Modesty, Humility, Truthfulness, Contentment, Tazkiyah/Ikhlas \\
\hline
3. Branches & The Family Unit & 13–18 & Marriage, Affection/Mercy, Simplicity, Parents, Children, Pedagogy \\
\hline
4. Leaves & Living Islam in the World & 19–25 & Sabr, Wasatiyyah, Joy, Beauty, Cleanliness, Time, Boundaries \\
\hline
5. Fruit & Trust, Hope, Eternal Vision & 26–30 & Justice, Charity, Inner Struggle, Dawah, Surrender \\
\hline
\end{longtable}

\textbf{Total: 5 phases - 30 modules - 30 weeks:} 

\vspace{6pt}\noindent{\color{border}\rule{\textwidth}{0.5pt}}\vspace{6pt}

\subsection*{Source Citation Index}

\subsubsection*{Books Used (all by Yusuf al-Qaradawi)}

\begin{longtable}{|p{2cm}|p{5cm}|p{8.5cm}|}
\hline
\textbf{\color{accent}\#} & \textbf{\color{accent}Book} & \textbf{\color{accent}Chapters/Sections Used} \\
\hline
1 & Ethics in Islam (أخلاق المسلم) & ch-01, ch-02, ch-03, ch-04, sections 04-01, 04-02, 04-03 \\
\hline
2 & Faith and Life (الإيمان والحياة) & ch-01, ch-02, ch-03, ch-04 \\
\hline
3 & The Lawful and the Prohibited in Islam & ch-03 (Marriage and Family Life) \\
\hline
4 & Diversion and Arts in Islam & ch-01, ch-02, ch-04, ch-07 \\
\hline
5 & Education and Economy in the Sunnah & ch-01 (Education) \\
\hline
6 & Economic Security in Islam & ch-01, ch-04 \\
\hline
7 & Approaching the Sunnah & ch-01 \\
\hline
8 & Fiqh al-Zakah & ch-01 (referenced for marriage/zakat intersection) \\
\hline
\end{longtable}

\subsubsection*{Qur'anic Verses Cited (by module)}

\begin{longtable}{|p{4.5cm}|p{11cm}|}
\hline
\textbf{\color{accent}Module} & \textbf{\color{accent}Key Verses} \\
\hline
1.1 & 17:70, 95:4, 49:13, 3:110 \\
\hline
1.2 & 13:28, 12:87 \\
\hline
1.3 & 30:21, 67:3 \\
\hline
1.4 & 12:87, 15:56, 65:2-3 \\
\hline
1.5 & 49:13, 2:183, 2:197, 5:2 \\
\hline
1.6 & 65:3, 3:160, 4:81 \\
\hline
2.1 & 68:4, 2:177 \\
\hline
2.2 & 24:30-31 \\
\hline
2.3 & 17:24, 17:37, 31:18, 63:8, 5:54 \\
\hline
2.4 & 33:35 \\
\hline
2.5 & 2:143, 25:67, 7:31, 57:23 \\
\hline
2.6 & 91:9-10, 98:5, 18:110 \\
\hline
3.1 & 30:21, 4:21, 24:32, 4:4, 25:54 \\
\hline
3.2 & 30:21, 4:19, 2:228 \\
\hline
3.3 & 4:4, 24:33, 65:7 \\
\hline
3.4 & 17:23-24, 46:15, 31:14-15 \\
\hline
3.5 & 17:31, 6:152, 4:10 \\
\hline
3.6 & 66:6 \\
\hline
4.1 & 2:153, 91:9-10, 29:69 \\
\hline
4.2 & 2:143, 25:67 \\
\hline
4.3 & 90:4, 7:31-32 \\
\hline
4.4 & 16:5-6, 16:8, 16:14, 6:141, 7:31-32 \\
\hline
4.5 & 2:222, 74:4, 7:31 \\
\hline
4.6 & 103:1-3, 63:10-11, 16:12 \\
\hline
4.7 & 17:32, 24:30-31 \\
\hline
5.1 & 16:90, 4:135, 5:8, 4:58 \\
\hline
5.2 & 6:164, 33:72, 51:19, 2:215 \\
\hline
5.3 & 91:7-10, 79:40-41, 29:69, 38:26 \\
\hline
5.4 & 16:125, 33:36 \\
\hline
5.5 & 65:2-3, 3:159, 2:156-157 \\
\hline
\end{longtable}

\vspace{6pt}\noindent{\color{border}\rule{\textwidth}{0.5pt}}\vspace{6pt}

\subsection*{Guidelines for the Halaqah Facilitator}

\subsubsection*{Before Each Session}
\begin{enumerate}
\item Read the relevant Qaradawi source text (cited in each module)
\item Prepare the Qur'an recitation (learn or find audio recitation)
\item Review the discussion questions and adapt them to your family's ages
\item Prepare the activity materials if needed
\end{enumerate}


\subsubsection*{During the Session}
\begin{enumerate}
\item Start on time, with du'a --- this sets the tone
\item Let children participate at their level --- they can recite Qur'an, share in discussion, do activities
\item Use the Prophetic pedagogy: gradual, merciful, questions, stories, moderate length
\item Don't lecture for 15 minutes straight --- break up the teaching with questions
\item End with the du'a --- let everyone say it together
\end{enumerate}


\subsubsection*{Adapting for Age}
\begin{itemize}
\item \textbf{Under 7:} Focus on the story, the du'a, and the activity. Simplify discussion to one question.
\item \textbf{7–12:} Full participation. Use the stories (Prophets, Companions) as anchors.
\item \textbf{13–17:} Full participation. Encourage them to challenge, question, and apply to their own lives.
\item \textbf{Adults:} Full depth. Read the Qaradawi source text beforehand for richer discussion.
\end{itemize}


\subsubsection*{If You Miss a Week}
\begin{itemize}
\item Don't skip the module --- just continue next week. The course is 30 weeks but can stretch to 35–40 with breaks for Eid, travel, exams, etc.
\item Hermes cron does NOT backfill missed runs --- and neither should you. Just pick up where you left off.
\end{itemize}


\vspace{6pt}\noindent{\color{border}\rule{\textwidth}{0.5pt}}\vspace{6pt}

\textit{This curriculum was compiled from the Qaradawi Library LLM Wiki at /root/qaradawi-library/. Every teaching point traces to the original works of Shaykh Yusuf al-Qaradawi. Qur'anic verses and hadiths are cited within each module. This document is for family educational use.}

\end{document}
